% Options for packages loaded elsewhere
\PassOptionsToPackage{unicode}{hyperref}
\PassOptionsToPackage{hyphens}{url}
\documentclass[
]{article}
\usepackage{xcolor}
\usepackage{amsmath,amssymb}
\setcounter{secnumdepth}{-\maxdimen} % remove section numbering
\usepackage{iftex}
\ifPDFTeX
  \usepackage[T1]{fontenc}
  \usepackage[utf8]{inputenc}
  \usepackage{textcomp} % provide euro and other symbols
\else % if luatex or xetex
  \usepackage{unicode-math} % this also loads fontspec
  \defaultfontfeatures{Scale=MatchLowercase}
  \defaultfontfeatures[\rmfamily]{Ligatures=TeX,Scale=1}
\fi
\usepackage{lmodern}
\ifPDFTeX\else
  % xetex/luatex font selection
\fi
% Use upquote if available, for straight quotes in verbatim environments
\IfFileExists{upquote.sty}{\usepackage{upquote}}{}
\IfFileExists{microtype.sty}{% use microtype if available
  \usepackage[]{microtype}
  \UseMicrotypeSet[protrusion]{basicmath} % disable protrusion for tt fonts
}{}
\makeatletter
\@ifundefined{KOMAClassName}{% if non-KOMA class
  \IfFileExists{parskip.sty}{%
    \usepackage{parskip}
  }{% else
    \setlength{\parindent}{0pt}
    \setlength{\parskip}{6pt plus 2pt minus 1pt}}
}{% if KOMA class
  \KOMAoptions{parskip=half}}
\makeatother
\usepackage{color}
\usepackage{fancyvrb}
\newcommand{\VerbBar}{|}
\newcommand{\VERB}{\Verb[commandchars=\\\{\}]}
\DefineVerbatimEnvironment{Highlighting}{Verbatim}{commandchars=\\\{\}}
% Add ',fontsize=\small' for more characters per line
\newenvironment{Shaded}{}{}
\newcommand{\AlertTok}[1]{\textcolor[rgb]{1.00,0.00,0.00}{\textbf{#1}}}
\newcommand{\AnnotationTok}[1]{\textcolor[rgb]{0.38,0.63,0.69}{\textbf{\textit{#1}}}}
\newcommand{\AttributeTok}[1]{\textcolor[rgb]{0.49,0.56,0.16}{#1}}
\newcommand{\BaseNTok}[1]{\textcolor[rgb]{0.25,0.63,0.44}{#1}}
\newcommand{\BuiltInTok}[1]{\textcolor[rgb]{0.00,0.50,0.00}{#1}}
\newcommand{\CharTok}[1]{\textcolor[rgb]{0.25,0.44,0.63}{#1}}
\newcommand{\CommentTok}[1]{\textcolor[rgb]{0.38,0.63,0.69}{\textit{#1}}}
\newcommand{\CommentVarTok}[1]{\textcolor[rgb]{0.38,0.63,0.69}{\textbf{\textit{#1}}}}
\newcommand{\ConstantTok}[1]{\textcolor[rgb]{0.53,0.00,0.00}{#1}}
\newcommand{\ControlFlowTok}[1]{\textcolor[rgb]{0.00,0.44,0.13}{\textbf{#1}}}
\newcommand{\DataTypeTok}[1]{\textcolor[rgb]{0.56,0.13,0.00}{#1}}
\newcommand{\DecValTok}[1]{\textcolor[rgb]{0.25,0.63,0.44}{#1}}
\newcommand{\DocumentationTok}[1]{\textcolor[rgb]{0.73,0.13,0.13}{\textit{#1}}}
\newcommand{\ErrorTok}[1]{\textcolor[rgb]{1.00,0.00,0.00}{\textbf{#1}}}
\newcommand{\ExtensionTok}[1]{#1}
\newcommand{\FloatTok}[1]{\textcolor[rgb]{0.25,0.63,0.44}{#1}}
\newcommand{\FunctionTok}[1]{\textcolor[rgb]{0.02,0.16,0.49}{#1}}
\newcommand{\ImportTok}[1]{\textcolor[rgb]{0.00,0.50,0.00}{\textbf{#1}}}
\newcommand{\InformationTok}[1]{\textcolor[rgb]{0.38,0.63,0.69}{\textbf{\textit{#1}}}}
\newcommand{\KeywordTok}[1]{\textcolor[rgb]{0.00,0.44,0.13}{\textbf{#1}}}
\newcommand{\NormalTok}[1]{#1}
\newcommand{\OperatorTok}[1]{\textcolor[rgb]{0.40,0.40,0.40}{#1}}
\newcommand{\OtherTok}[1]{\textcolor[rgb]{0.00,0.44,0.13}{#1}}
\newcommand{\PreprocessorTok}[1]{\textcolor[rgb]{0.74,0.48,0.00}{#1}}
\newcommand{\RegionMarkerTok}[1]{#1}
\newcommand{\SpecialCharTok}[1]{\textcolor[rgb]{0.25,0.44,0.63}{#1}}
\newcommand{\SpecialStringTok}[1]{\textcolor[rgb]{0.73,0.40,0.53}{#1}}
\newcommand{\StringTok}[1]{\textcolor[rgb]{0.25,0.44,0.63}{#1}}
\newcommand{\VariableTok}[1]{\textcolor[rgb]{0.10,0.09,0.49}{#1}}
\newcommand{\VerbatimStringTok}[1]{\textcolor[rgb]{0.25,0.44,0.63}{#1}}
\newcommand{\WarningTok}[1]{\textcolor[rgb]{0.38,0.63,0.69}{\textbf{\textit{#1}}}}
\usepackage{longtable,booktabs,array}
\newcounter{none} % for unnumbered tables
\usepackage{calc} % for calculating minipage widths
% Correct order of tables after \paragraph or \subparagraph
\usepackage{etoolbox}
\makeatletter
\patchcmd\longtable{\par}{\if@noskipsec\mbox{}\fi\par}{}{}
\makeatother
% Allow footnotes in longtable head/foot
\IfFileExists{footnotehyper.sty}{\usepackage{footnotehyper}}{\usepackage{footnote}}
\makesavenoteenv{longtable}
\setlength{\emergencystretch}{3em} % prevent overfull lines
\providecommand{\tightlist}{%
  \setlength{\itemsep}{0pt}\setlength{\parskip}{0pt}}
\usepackage{bookmark}
\IfFileExists{xurl.sty}{\usepackage{xurl}}{} % add URL line breaks if available
\urlstyle{same}
\hypersetup{
  pdftitle={CAIRN Protocol Whitepaper},
  pdfauthor={Maroua Boudoukha},
  hidelinks,
  pdfcreator={LaTeX via pandoc}}

\title{CAIRN Protocol Whitepaper}
\author{Maroua Boudoukha}
\date{May 2026}

\begin{document}
\maketitle

{
\setcounter{tocdepth}{3}
\tableofcontents
}
\section{Agent Failure and Recovery
Protocol}\label{agent-failure-and-recovery-protocol}

\subsection{Version 2.0 --- April 2026}\label{version-2.0-april-2026}

\begin{quote}
\textbf{Author:} Maroua Boudoukha \textbf{Affiliation:} Independent
Researcher \textbf{Contact:} github.com/MarouaBoud \textbf{ORCID:}
\emph{to be added before formal submission} \textbf{Email:} \emph{to be
added before formal submission}

\textbf{Copyright 2026 Maroua BOUDOUKHA. All rights reserved.}

This document may be cited for academic and research purposes with
proper attribution: BOUDOUKHA, M. (2026). \emph{CAIRN Protocol: Agent
Failure and Recovery Protocol}. Whitepaper v2.0, April 2026.

Redistribution or commercial use requires written permission from the
author.
\end{quote}

\begin{center}\rule{0.5\linewidth}{0.5pt}\end{center}

\section{Abstract}\label{abstract}

AI agent task completion rates remain at approximately 50\% across
popular frameworks, yet no standardized protocol exists for failure
detection, classification, and recovery in the on-chain agent economy.
We present CAIRN, the \textbf{first on-chain agent protocol} to classify
agent failures by \textbf{recoverability} rather than symptom ---
adapting the classical crash-vs-Byzantine distinction from distributed
systems {[}8{]} to the AI agent domain --- enabling deterministic
routing to checkpoint-based recovery or dispute resolution.

CAIRN defines a 6-state machine with three-tier recovery routing,
enforced by smart contracts: when an agent fails mid-task, the protocol
detects the failure via missed heartbeats or resource exhaustion,
classifies it into one of three recoverability classes (LIVENESS,
RESOURCE, LOGIC), computes a multiplicative recovery score \emph{r} =
\emph{F}\textsuperscript{0.80} ×\textbackslash{}
\emph{B}\textsuperscript{0.35} ×\textbackslash{}
\emph{D}\textsuperscript{0.15}, and routes the task to either a
qualified fallback agent who resumes from the last IPFS-committed
checkpoint, or to dispute resolution. The formula is \textbf{calibrated
and validated via Monte Carlo simulation} against a ground-truth model
derived from published failure-mode distributions {[}1{]}{[}2{]}{[}3{]}:
across 100,000 synthetic task-failure events and 16 experiments, the
multiplicative formula achieves 23.46\% misrouting against that ground
truth --- within 0.93pp of the Bayes-optimal minimum (22.53\%)
attainable for the same model --- and reduces wasted-recovery false
positives by 65\% versus a linear baseline. We are explicit that this is
near-optimality \emph{against the calibrated model}, not against
measured production data, which does not yet exist; the staged
calibration roadmap (Section 10.1) replaces the synthetic ground truth
with observed outcomes as testnet and mainnet data accumulate. Escrow is
settled proportionally to verified work. We prove escrow safety,
termination, and state determinism, and show that honest checkpointing
is the dominant strategy under realistic economic parameters.

Our key insight is that \textbf{economic enforcement} ---
escrow-conditioned record writing --- bootstraps a collective
intelligence layer without requiring altruistic participation. Every
failure becomes a queryable record. Every recovery teaches the next
agent. The accumulated execution history is on-chain and publicly
indexable, so the \emph{records} themselves are not forkable in the
trivial sense, but replicating CAIRN's \emph{useful} intelligence layer
requires reproducing the escrow-mandated write volume --- a competing
protocol would need to independently generate comparable failure data,
which presupposes comparable task throughput. This is a network effect
with a switching-cost moat, not a structural impossibility.

CAIRN integrates three Ethereum standards: ERC-8004 for agent identity
and reputation, ERC-8183 for job escrow lifecycle, and ERC-7710 for
scoped delegation. It is deployed on Base and composable with existing
agent frameworks (LangGraph, Olas, CrewAI, AutoGen) and emerging
coordination protocols (Google A2A, Anthropic MCP).

\begin{quote}
\textbf{Scope of empirical claims.} The 23.46\% misrouting result is a
simulation outcome against a literature-calibrated ground-truth model.
We acknowledge three concrete bounds on its external validity: (i) the
ground-truth recovery rates and class frequencies are \emph{assumed}
from prior published studies, not \emph{measured} from CAIRN's own
deployment; (ii) full checkpoint portability holds for
structured-pipeline tasks but degrades for reasoning-heavy workloads
(Section 4.1.1); (iii) the arbiter incentive analysis (Section 7.5
Proposition 3) assumes a detection probability \emph{p}\textsubscript{d}
that is high for LIVENESS/RESOURCE disputes but plausibly much lower for
LOGIC disputes where the underlying question is ``did the agent reason
incorrectly?'' Section 10.3 enumerates the full limitation set.

\textbf{Reproducibility.} All simulation code, results, and figures are
reproducible from \texttt{simulation/} in the CAIRN repository via
\texttt{python3\ -m\ simulation.run\_eq4} (seed=42, deterministic on
NumPy ≥1.20). See reference {[}18{]}.
\end{quote}

\begin{quote}
\textbf{Note on protocol versions.} This paper specifies CAIRN
\textbf{v2}, the simulation-validated protocol described throughout. The
current testnet deployment (\textbf{v1}) uses an interim linear recovery
score (Equation 1 of Section 10.1) with binary routing, reflecting the
protocol's state prior to the calibration work reported here. The
multiplicative formula, three-tier routing, and refined stake/threshold
parameters described in Sections 2.2.1, 6.4, and 7.5 are the v2
specification, intended for adoption through the governance upgrade path
outlined in Section 8.3. Sections marked \emph{``v2 specification''}
describe the target protocol; \emph{``v1 deployment''} references
describe what is live on testnet. This paper exists to motivate and
document the v1 → v2 transition.
\end{quote}

\begin{center}\rule{0.5\linewidth}{0.5pt}\end{center}

\section{Table of Contents}\label{table-of-contents}

\begin{enumerate}
\def\labelenumi{\arabic{enumi}.}
\tightlist
\item
  \hyperref[1-introduction]{Introduction}
\item
  \hyperref[2-protocol-specification]{Protocol Specification}
\item
  \hyperref[3-design-philosophy]{Design Philosophy}
\item
  \hyperref[4-key-design-decisions]{Key Design Decisions}
\item
  \hyperref[5-execution-intelligence-layer]{Execution Intelligence
  Layer}
\item
  \hyperref[6-economic-model]{Economic Model}
\item
  \hyperref[7-security-model]{Security Model}
\item
  \hyperref[8-governance]{Governance}
\item
  \hyperref[9-related-work]{Related Work}
\item
  \hyperref[10-future-work]{Future Work}
\item
  \hyperref[11-references]{References}
\end{enumerate}

\begin{center}\rule{0.5\linewidth}{0.5pt}\end{center}

\section{1. Introduction}\label{introduction}

\subsection{1.1 The Problem}\label{the-problem}

The Ethereum agentic economy generates significant economic activity ---
over 10 million agent-to-agent transactions on the Olas network alone
{[}5{]},\textbackslash{} \textasciitilde49,000 registered agent
identities via ERC-8004 across 30+ EVM chains (as of February 2026)
{[}14{]}, and growing on-chain commerce via ERC-8183 {[}15{]}. But every
agent is operationally isolated.

When an agent fails mid-task --- because an API rate-limits, a budget is
exceeded, a context window overflows, or a process crashes ---
\textbf{nothing standard happens}. The escrow sits in an ambiguous
state. The operator who submitted the task --- whether a human, another
agent, a DAO multisig, or an autonomous smart contract --- has no
standardized way to detect the failure or trigger a recovery; the
failure is discovered (if at all) only when the principal happens to
poll the task. Another agent does not automatically take over. Completed
work is lost.

Twenty minutes later, a different agent --- same task type, same API,
same conditions --- fails identically. The collective cost compounds as
the agent economy scales.

\subsection{1.2 The Evidence}\label{the-evidence}

Published research establishes agent failure as a systemic problem, not
an edge case:

\begin{itemize}
\item
  Multi-agent benchmarks show an \textbf{average task completion rate of
  approximately 50\%} across popular frameworks including TaskWeaver,
  MetaGPT, and AutoGen {[}3{]}. The compounding effect is structural:
  \emph{even if} single-step accuracy were 85\% (well above current
  measured agent reliability per {[}2{]}), a 10-step workflow would
  succeed only \textasciitilde20\% of the time\textbackslash{}
  (0.85\textsuperscript{10} = 0.197). Real production accuracy is lower,
  and the 50\% headline rate is the empirical consequence.
\item
  The MAST taxonomy identifies \textbf{14 distinct failure modes} across
  1,600+ annotated traces from 7 multi-agent frameworks {[}1{]}.
  However, MAST classifies failures by symptom (step repetition,
  incorrect tool selection) --- not by what recovery action to take.
\item
  Research on AI agent reliability finds that \textbf{consistency
  remains weak across all models} --- outcome consistency is the most
  persistently low dimension, and agents cannot reliably determine when
  they are wrong {[}2{]}. This validates the need for external failure
  detection infrastructure rather than relying on agent self-diagnosis.
\item
  A systematic survey of 317 publications on autonomous agents and
  blockchains identifies \textbf{missing interface layers and verifiable
  policy enforcement} as key gaps {[}6{]} --- precisely the gaps CAIRN
  addresses.
\item
  ISO/IEC TR 5469 notes that ``highly autonomous systems likely fall
  into a risk category where current methods are insufficient to
  adequately mitigate reliability-related risks'' {[}7{]}.
\end{itemize}

\subsection{1.3 What Is Missing}\label{what-is-missing}

Every team building agents has written bespoke, incompatible failure
handling. There is:

\begin{itemize}
\tightlist
\item
  No standard definition of what an agent failure is
\item
  No standard protocol for what happens when one is detected
\item
  No standard mechanism for task handoff to a fallback agent
\item
  No standard escrow settlement rule for partial completion
\item
  No shared record of what failed, why, and what worked instead
\end{itemize}

CAIRN fills this gap.

\subsection{1.4 Scope and Limitations}\label{scope-and-limitations}

Before describing the protocol in detail, we surface the bounds within
which its claims hold. These are expanded in Section 10.3.

\begin{itemize}
\tightlist
\item
  \textbf{Empirical scope.} The 23.46\% misrouting result is from Monte
  Carlo simulation against a ground-truth model calibrated to published
  failure-mode distributions {[}1{]}{[}2{]}{[}3{]}. Class frequencies
  and base recovery rates are \emph{assumed} from prior literature, not
  \emph{measured} from CAIRN's own deployment, which does not yet exist
  at scale. The staged calibration roadmap (Section 10.1) replaces the
  synthetic model with observed outcomes as testnet and mainnet data
  accumulate.
\item
  \textbf{Task scope.} Full checkpoint portability covers
  structured-pipeline tasks (data fetches, API calls, multi-step
  computations) and tasks whose state can be made portable by
  serialising explicit context. For reasoning-heavy workloads with
  implicit framework state --- chain-of-thought, multi-turn dialogue,
  planning with backtracking --- only output-level checkpoints are
  portable and the recovery quality degrades (Section 4.1.1).
\item
  \textbf{Arbiter scope.} The economic analysis of arbiter honesty
  (Proposition 3 in Section 7.5) holds strongly for LIVENESS and
  RESOURCE disputes whose evidence is fully on-chain. For LOGIC disputes
  (``did the agent hallucinate?'', ``did the output match the
  specification?'') the detection probability is materially lower; CAIRN
  inherits the structural limitation of all on-chain arbitration
  protocols {[}19{]}. The v2 specification ships single-tier
  arbitration; an appeals layer is reserved for v3.
\item
  \textbf{Trust scope.} The protocol assumes the underlying L2 sequencer
  (Base/Coinbase, as of 2026) is available and does not actively collude
  with a dispute party. This is weaker than the L1 trust model and is
  the principal residual trust dependency in v2 (Section 7.2).
\item
  \textbf{Implementation scope.} This paper specifies the v2 protocol;
  the testnet currently runs v1 (interim linear formula, binary routing,
  15\% arbiter stake, no schema-hash enforcement). The v2 migration path
  is governance-gated through the \texttt{IRecoveryRouter} interface
  (Section 8.3) and a Phase-1 reference implementation of the v2 router
  with measured gas is included in the companion repository.
\end{itemize}

\begin{center}\rule{0.5\linewidth}{0.5pt}\end{center}

\section{2. Protocol Specification}\label{protocol-specification}

\subsection{2.1 Overview}\label{overview}

CAIRN is a standardized agent failure and recovery protocol. It defines
the exact sequence of events that occur when an agent fails mid-task ---
from detection, through classification, through fallback assignment,
through settlement --- \textbf{without any human-in-the-loop after task
submission}, and without requiring trust between agents.

\textbf{Definition (Operator).} Throughout this paper, an
\emph{operator} is the Ethereum address that submits a task and posts
its escrow --- an EOA, an agent, a smart contract, or a DAO multisig;
CAIRN is agnostic to which. The ``no human required'' property is
therefore precise: \textbf{no human signature, intervention, or polling
is required at any point between task submission and final settlement}.
Failure detection, classification, recovery routing, fallback execution,
dispute initiation, and escrow distribution all proceed via
permissionless on-chain enforcement (Section 3.3) regardless of who or
what submitted the task.

As a byproduct, CAIRN accumulates an \textbf{execution intelligence
layer}: a shared, queryable record of every failure, recovery, and
completion across the ecosystem. The intelligence layer grows
automatically because record-writing is mandatory for escrow settlement.

\textbf{Formal Definition.} A CAIRN task is a tuple \emph{T} =
(\emph{id}, \emph{σ},\textbackslash{}
\emph{A}\textsubscript{p},\textbackslash{} \emph{A}\textsubscript{f},
\emph{E}, \emph{δ}, \emph{H},\textbackslash{}
\emph{c}\textsubscript{p},\textbackslash{} \emph{c}\textsubscript{f},
\emph{κ},\textbackslash{} \emph{t}\textsubscript{0}) where:

{\def\LTcaptype{none} % do not increment counter
\begin{longtable}[]{@{}
  >{\raggedright\arraybackslash}p{(\linewidth - 4\tabcolsep) * \real{0.2963}}
  >{\raggedright\arraybackslash}p{(\linewidth - 4\tabcolsep) * \real{0.2222}}
  >{\raggedright\arraybackslash}p{(\linewidth - 4\tabcolsep) * \real{0.4815}}@{}}
\toprule\noalign{}
\begin{minipage}[b]{\linewidth}\raggedright
Symbol
\end{minipage} & \begin{minipage}[b]{\linewidth}\raggedright
Type
\end{minipage} & \begin{minipage}[b]{\linewidth}\raggedright
Description
\end{minipage} \\
\midrule\noalign{}
\endhead
\bottomrule\noalign{}
\endlastfoot
\emph{id} & bytes32 & Unique task identifier: keccak256(operator, nonce,
block.timestamp) \\
\emph{σ} & enum & Current state ∈ \{IDLE, RUNNING, FAILED, RECOVERING,
DISPUTED, RESOLVED\} \\
\textbackslash{} \emph{A}\textsubscript{p},\textbackslash{}
\emph{A}\textsubscript{f} & address & Primary and fallback agent
addresses \\
\emph{E} & uint256 & Escrow amount in wei, \emph{E} ≥\textbackslash{}
\emph{E}\textsubscript{min} =\textbackslash{} 10\textsuperscript{15}
(0.001 ETH) \\
\emph{δ} & uint256 & Deadline as block timestamp \\
\emph{H} & uint256 & Heartbeat interval in seconds, 30 ≤ \emph{H} ≤
\emph{δ}/4 \\
\textbackslash{} \emph{c}\textsubscript{p},\textbackslash{}
\emph{c}\textsubscript{f} & uint256 & Checkpoint counts for primary and
fallback agents \\
\emph{κ} & uint256 & Cost accrued in wei \\
\textbackslash{} \emph{t}\textsubscript{0} & uint256 & Task start block
timestamp \\
\end{longtable}
}

\subsection{2.2 State Machine}\label{state-machine}

Six states. Every transition is deterministic. After the operator
submits the task, \textbf{no human signature or human intervention is
required to trigger any state change} --- every transition fires from
on-chain conditions evaluated by permissionless enforcement functions
(Section 3.3) that any address may call.

\begin{verbatim}
                    ┌──────────────────────────────────────────────────┐
                    │               CAIRN State Machine                │
                    └──────────────────────────────────────────────────┘

    ┌──────┐  confirm   ┌─────────┐   task done   ┌──────────┐
    │      │ ─────────► │         │ ─────────────► │          │
    │ IDLE │            │ RUNNING │                │ RESOLVED │ (terminal)
    │      │            │         │                │          │
    └──────┘            └────┬────┘                └────▲─────┘
                             │                          │
                    fault    │                          │ complete
                  detected   │                          │
                             ▼                          │
                        ┌─────────┐  score ≥ 0.35  ┌───┴──────┐
                        │         │ ──────────────► │          │
                        │ FAILED  │  (≥0.40 full   │RECOVERING│
                        │         │ ◄─0.35-0.40─── │ (full or │
                        └────┬────┘  reduced scope) │ reduced) │
                             │       or fallback    └──────────┘
                      score  │       fails
                     < 0.35  │
                             ▼
                        ┌──────────┐  arbiter   ┌──────────┐
                        │          │ ──────────► │          │
                        │ DISPUTED │             │ RESOLVED │ (terminal)
                        │          │  timeout    │          │
                        └──────────┘ ──────────► └──────────┘
                                      refund
\end{verbatim}

\textbf{State Definitions:}

{\def\LTcaptype{none} % do not increment counter
\begin{longtable}[]{@{}
  >{\raggedright\arraybackslash}p{(\linewidth - 6\tabcolsep) * \real{0.1458}}
  >{\raggedright\arraybackslash}p{(\linewidth - 6\tabcolsep) * \real{0.3125}}
  >{\raggedright\arraybackslash}p{(\linewidth - 6\tabcolsep) * \real{0.1875}}
  >{\raggedright\arraybackslash}p{(\linewidth - 6\tabcolsep) * \real{0.3542}}@{}}
\toprule\noalign{}
\begin{minipage}[b]{\linewidth}\raggedright
State
\end{minipage} & \begin{minipage}[b]{\linewidth}\raggedright
Entry Trigger
\end{minipage} & \begin{minipage}[b]{\linewidth}\raggedright
Actions
\end{minipage} & \begin{minipage}[b]{\linewidth}\raggedright
Exit Conditions
\end{minipage} \\
\midrule\noalign{}
\endhead
\bottomrule\noalign{}
\endlastfoot
\textbf{IDLE} & Operator submits task & Query intelligence layer; lock
escrow; pre-authorize delegation & Operator confirms → RUNNING \\
\textbf{RUNNING} & Operator confirmation & Agent executes subtasks;
writes checkpoints; emits heartbeats & Success → RESOLVED; Fault →
FAILED \\
\textbf{FAILED} & Liveness miss, budget hit, or deadline exceeded &
Classify failure; compute recovery score; write failure record & Score ≥
0.35 → RECOVERING (full or reduced scope); Score \textless{} 0.35 →
DISPUTED \\
\textbf{RECOVERING} & Recovery score ≥ 0.35 & Select fallback; transfer
checkpoint state; fallback resumes (full or reduced scope) & Success →
RESOLVED; Failure → DISPUTED \\
\textbf{DISPUTED} & Score \textless{} 0.35 or fallback failure & Hold
escrow; expose evidence; start arbiter timeout & Arbiter ruling →
RESOLVED; Timeout → auto-refund \\
\textbf{RESOLVED} & Completion or ruling & Settle escrow proportionally;
write resolution record; update reputation & Terminal \\
\end{longtable}
}

\subsection{2.2.1 Formal Properties}\label{formal-properties}

\textbf{Transition Function.} The state transition function \emph{τ}:
\emph{S} × \emph{Event} → \emph{S} is defined as:

{\def\LTcaptype{none} % do not increment counter
\begin{longtable}[]{@{}
  >{\raggedright\arraybackslash}p{(\linewidth - 6\tabcolsep) * \real{0.3333}}
  >{\raggedright\arraybackslash}p{(\linewidth - 6\tabcolsep) * \real{0.1556}}
  >{\raggedright\arraybackslash}p{(\linewidth - 6\tabcolsep) * \real{0.2444}}
  >{\raggedright\arraybackslash}p{(\linewidth - 6\tabcolsep) * \real{0.2667}}@{}}
\toprule\noalign{}
\begin{minipage}[b]{\linewidth}\raggedright
Current State
\end{minipage} & \begin{minipage}[b]{\linewidth}\raggedright
Event
\end{minipage} & \begin{minipage}[b]{\linewidth}\raggedright
Condition
\end{minipage} & \begin{minipage}[b]{\linewidth}\raggedright
Next State
\end{minipage} \\
\midrule\noalign{}
\endhead
\bottomrule\noalign{}
\endlastfoot
IDLE & Confirm & Operator address calls \texttt{confirmTask} (signs via
wallet, contract call, or multisig) & RUNNING \\
RUNNING & Complete & All subtasks verified & RESOLVED \\
RUNNING & HeartbeatMiss & block.timestamp \textgreater{} lastHeartbeat +
\emph{H} & FAILED \\
RUNNING & BudgetExceeded & \emph{κ} ≥ \emph{E} & FAILED \\
RUNNING & DeadlineExceeded & block.timestamp ≥ \emph{δ} & FAILED \\
FAILED & RecoveryRoute & \emph{r} ≥ 0.35 & RECOVERING \\
FAILED & RecoveryRoute & \emph{r} \textless{} 0.35 & DISPUTED \\
RECOVERING & Complete & Fallback completes remaining subtasks &
RESOLVED \\
RECOVERING & FallbackFailed & Fallback fails or deadline exceeded &
DISPUTED \\
DISPUTED & ArbiterRuling & Arbiter submits valid ruling & RESOLVED \\
DISPUTED & Timeout & block.timestamp ≥\textbackslash{}
\emph{t}\textsubscript{dispute} +\textbackslash{}
\emph{D}\textsubscript{timeout} & RESOLVED (refund) \\
\end{longtable}
}

All (\emph{σ}, \emph{event}) pairs not listed above are undefined; the
transaction reverts.

\textbf{Theorem 1 (Escrow Safety).} \emph{For any task T, escrow E is
not distributed until σ = RESOLVED.}

\emph{Proof.} The settlement function \texttt{settle(taskId)} contains
the precondition \texttt{require(task.state\ ==\ RESOLVED)}. No other
function in the protocol transfers escrow from the task's balance. The
transition table in Section 2.2 has no outgoing edges from RESOLVED, so
RESOLVED is terminal. Therefore, escrow \emph{E} remains locked in all
non-terminal states and is released only via \texttt{settle()} at σ =
RESOLVED. ∎

\textbf{Theorem 2 (Termination).} \emph{Every task T reaches σ =
RESOLVED within at most δ − t₀ +\textbackslash{}
D\textsubscript{timeout} seconds.}

\emph{Proof.} We show each non-terminal state has a bounded exit: -
\textbf{IDLE}: Operator confirms or deadline passes. Bounded by \emph{δ}
−\textbackslash{} \emph{t}\textsubscript{0}. - \textbf{RUNNING}: Either
completes (→ RESOLVED) or a fault triggers (heartbeat miss at \emph{t}
\textgreater{} lastHeartbeat + \emph{H}, or deadline \emph{δ} reached).
Maximum duration: \emph{δ} −\textbackslash{} \emph{t}\textsubscript{0}.
- \textbf{FAILED}: Recovery score \emph{r} is computed as a pure
function of on-chain state. Routing is immediate within the same
transaction. Duration: 0. - \textbf{RECOVERING}: Fallback either
completes (→ RESOLVED) or fails (→ DISPUTED). Bounded by remaining
deadline: \emph{δ} −\textbackslash{} \emph{t}\textsubscript{current}. -
\textbf{DISPUTED}: Arbiter rules (→ RESOLVED) or timeout expires (→
RESOLVED with auto-refund). Bounded by\textbackslash{}
\emph{D}\textsubscript{timeout} (default 604,800 seconds = 7 days).

RUNNING and RECOVERING share the same time window\textbackslash{}
{[}\emph{t}\textsubscript{0}, \emph{δ}{]} --- RECOVERING begins
at\textbackslash{} \emph{t}\textsubscript{fail} ≥\textbackslash{}
\emph{t}\textsubscript{0} and can only consume time up to \emph{δ}, so
the two intervals are not additive. The worst-case path therefore takes
at most (\emph{δ} −\textbackslash{} \emph{t}\textsubscript{0}) seconds
before DISPUTED is entered, plus at most\textbackslash{}
\emph{D}\textsubscript{timeout} seconds in DISPUTED. Total upper bound:
(\emph{δ} −\textbackslash{} \emph{t}\textsubscript{0}) +\textbackslash{}
\emph{D}\textsubscript{timeout}. ∎

\textbf{Theorem 3 (Irreversibility).} \emph{State transitions are
monotonic: once a task leaves state σ, it never returns to σ.}

\emph{Proof.} The transition function \emph{τ} defines a directed
acyclic graph over states: IDLE → RUNNING → \{FAILED, RESOLVED\}; FAILED
→ \{RECOVERING, DISPUTED\}; RECOVERING → \{RESOLVED, DISPUTED\};
DISPUTED → RESOLVED. No edge returns to a previously visited state. Each
transition executes atomically within a single transaction, writing the
new state to contract storage before function return. ∎

\textbf{Theorem 4 (Determinism).} \emph{For any task T in state σ and
event e, τ(σ, e) produces at most one successor state.}

\emph{Proof.} The only branching transition is \emph{τ}(FAILED,
RecoveryRoute(\emph{r})), which depends on the recovery score \emph{r}.
Since \emph{r} is computed as a pure function of on-chain state
(Equation 1 in Section 6.4), and the v2 threshold comparisons (\emph{r}
≥ 0.40, 0.35 ≤ \emph{r} \textless{} 0.40, \emph{r} \textless{} 0.35)
partition ℝ into disjoint intervals, exactly one branch is taken. (In
v1, the two-tier partition at \emph{r} ≥ 0.30 / \emph{r} \textless{}
0.30 is a degenerate case of the same proof; the determinism property
holds in both versions.) All other transitions in the table map to a
unique successor. ∎

\subsection{2.3 Architecture}\label{architecture}

\begin{verbatim}
┌─────────────────────────────────────────────────────────────────────┐
│                        CAIRN Protocol Architecture                  │
├─────────────────────────────────────────────────────────────────────┤
│                                                                     │
│   Operator / Agent Frameworks                                       │
│   (LangGraph, Olas SDK, CrewAI, AutoGen, custom)                    │
│         │                                                           │
│         ▼                                                           │
│   ┌───────────────────────────────────────────────────────────┐     │
│   │                    CairnCore (Hub)                         │     │
│   │  Task lifecycle · Checkpoint batching · Heartbeat          │     │
│   │  Escrow settlement · Permissionless enforcement            │     │
│   ├───────────┬───────────┬───────────┬──────────────────────┤     │
│   │ Recovery  │ Fallback  │ Arbiter   │ Governance            │     │
│   │ Router    │ Pool      │ Registry  │                       │     │
│   │           │           │           │ Timelock proposals     │     │
│   │ Classify  │ Select    │ Dispute   │ Parameter control      │     │
│   │ Score     │ Rank      │ Ruling    │ Upgrade authorization  │     │
│   │ Route     │ Slash     │ Appeal    │                       │     │
│   └─────┬─────┴─────┬─────┴─────┬─────┴──────────────────────┘     │
│         │           │           │                                   │
│   ┌─────▼─────┐ ┌───▼──────┐ ┌─▼───────────┐                      │
│   │ ERC-8183  │ │ ERC-8004 │ │  ERC-7710   │                      │
│   │ Job       │ │ Identity │ │  Delegation  │                      │
│   │ Escrow    │ │ + Reputa-│ │  (scoped    │                      │
│   │ Hooks     │ │ tion     │ │  permission) │                      │
│   └───────────┘ └──────────┘ └─────────────┘                      │
│         │           │                                               │
│   ┌─────▼───────────▼──────────────────────────────────────────┐   │
│   │              Execution Intelligence Layer                   │   │
│   │  IPFS (checkpoint + record storage)                         │   │
│   │  The Graph (event indexing + aggregation)                   │   │
│   │  Knowledge graph (pattern detection + queries)              │   │
│   └─────────────────────────────────────────────────────────────┘   │
│                                                                     │
│   Optional Integration:                                             │
│   ┌──────────────┐  ┌──────────────┐                               │
│   │ Olas Mech    │  │ OlasMech     │                               │
│   │ Marketplace  │  │ Adapter      │                               │
│   │\\ (~2000 mechs)│  │ (fallback    │                               │
│   │              │  │  bridge)     │                               │
│   └──────────────┘  └──────────────┘                               │
└─────────────────────────────────────────────────────────────────────┘
\end{verbatim}

\subsection{2.4 Worked Example: The 2:47am
Recovery}\label{worked-example-the-247am-recovery}

A DeFi portfolio-management agent (acting as the operator --- its own
wallet posts the escrow on behalf of an end user it serves) submits a
5-step portfolio rebalancing task to a worker agent with 0.01 ETH
escrow, a 30-minute deadline, and a 60-second heartbeat interval. The
operator role here is filled by an agent, not a human; the example would
proceed identically if the operator were a human user, a DAO treasury
contract, or any other principal.

{\def\LTcaptype{none} % do not increment counter
\begin{longtable}[]{@{}
  >{\raggedright\arraybackslash}p{(\linewidth - 4\tabcolsep) * \real{0.3000}}
  >{\raggedright\arraybackslash}p{(\linewidth - 4\tabcolsep) * \real{0.3500}}
  >{\raggedright\arraybackslash}p{(\linewidth - 4\tabcolsep) * \real{0.3500}}@{}}
\toprule\noalign{}
\begin{minipage}[b]{\linewidth}\raggedright
Time
\end{minipage} & \begin{minipage}[b]{\linewidth}\raggedright
Event
\end{minipage} & \begin{minipage}[b]{\linewidth}\raggedright
State
\end{minipage} \\
\midrule\noalign{}
\endhead
\bottomrule\noalign{}
\endlastfoot
T+0s & Operator submits task; escrow locked & \textbf{IDLE} \\
T+5s & Operator confirms; agent begins execution & \textbf{RUNNING} \\
T+30s & Agent completes step 1 → writes checkpoint CID to IPFS → commits
on-chain & RUNNING \\
T+65s & Agent completes step 2 → checkpoint 2 committed & RUNNING \\
T+95s & Agent completes step 3 → checkpoint 3 committed & RUNNING \\
T+120s & Agent calls CoinGecko API → HTTP 429 (rate limit) → agent
process crashes & RUNNING \\
T+185s & Heartbeat missed (120s + 65s \textgreater{} 60s interval) →
anyone calls \texttt{checkLiveness} & \textbf{FAILED} \\
T+186s & RecoveryRouter classifies: \textbf{LIVENESS} (\emph{F} = 0.70);
budget 85\% remaining; deadline 88\% remaining & FAILED \\
T+186s & Recovery score =\textbackslash{} \emph{F}\^{}0.80 ×
\emph{B}\^{}0.35 ×\textbackslash{} \emph{D}\^{}0.15 = 0.70\^{}0.80
×\textbackslash{} 0.85\^{}0.35 × 0.88\^{}0.15 = \textbf{0.752 × 0.945 ×
0.981 = 0.697} & FAILED \\
T+186s & Score 0.697 ≥ 0.40 → route to RECOVERING (full scope) &
\textbf{RECOVERING} \\
T+190s & FallbackPool selects highest-reputation agent for
\texttt{defi.trade\_execute} & RECOVERING \\
T+195s & Fallback reads checkpoints 1-3 from IPFS; resumes at step 4 &
RECOVERING \\
T+230s & Fallback completes step 4 → checkpoint 4 committed &
RECOVERING \\
T+270s & Fallback completes step 5 → task complete &
\textbf{RESOLVED} \\
T+271s & Settlement: primary gets 60\% (3/5 checkpoints), fallback gets
40\% (2/5), minus 0.5\% protocol fee & RESOLVED \\
\end{longtable}
}

\textbf{Result without CAIRN:} The operator (agent or human) discovers
the failure only on its next polling cycle --- typically 4+ hours later
for a human, or whenever the next health-check fires for an automated
principal. Full restart. Original agent paid 0.

\textbf{Result with CAIRN:} Automatic detection in 65 seconds. Fallback
resumes from step 4. Original agent paid 0.006 ETH for verified work
(60\% × 0.01 ETH escrow, minus 0.5\% protocol fee). Total recovery time:
\textasciitilde85 seconds.

\subsection{2.5 Comparative Analysis: Recovery
vs.~Restart}\label{comparative-analysis-recovery-vs.-restart}

The value of checkpoint-based recovery scales with task length and
failure point. The later a task fails, the more work is preserved:

{\def\LTcaptype{none} % do not increment counter
\begin{longtable}[]{@{}
  >{\raggedright\arraybackslash}p{(\linewidth - 16\tabcolsep) * \real{0.1111}}
  >{\raggedright\arraybackslash}p{(\linewidth - 16\tabcolsep) * \real{0.1111}}
  >{\raggedright\arraybackslash}p{(\linewidth - 16\tabcolsep) * \real{0.1111}}
  >{\raggedright\arraybackslash}p{(\linewidth - 16\tabcolsep) * \real{0.1111}}
  >{\raggedright\arraybackslash}p{(\linewidth - 16\tabcolsep) * \real{0.1111}}
  >{\raggedright\arraybackslash}p{(\linewidth - 16\tabcolsep) * \real{0.1111}}
  >{\raggedright\arraybackslash}p{(\linewidth - 16\tabcolsep) * \real{0.1111}}
  >{\raggedright\arraybackslash}p{(\linewidth - 16\tabcolsep) * \real{0.1111}}
  >{\raggedright\arraybackslash}p{(\linewidth - 16\tabcolsep) * \real{0.1111}}@{}}
\toprule\noalign{}
\begin{minipage}[b]{\linewidth}\raggedright
Task Profile
\end{minipage} & \begin{minipage}[b]{\linewidth}\raggedright
Steps
\end{minipage} & \begin{minipage}[b]{\linewidth}\raggedright
Failure Point
\end{minipage} & \begin{minipage}[b]{\linewidth}\raggedright
CAIRN Detection
\end{minipage} & \begin{minipage}[b]{\linewidth}\raggedright
Restart Detection
\end{minipage} & \begin{minipage}[b]{\linewidth}\raggedright
CAIRN Recovered Work
\end{minipage} & \begin{minipage}[b]{\linewidth}\raggedright
Restart Recovered Work
\end{minipage} & \begin{minipage}[b]{\linewidth}\raggedright
CAIRN Settlement Time
\end{minipage} & \begin{minipage}[b]{\linewidth}\raggedright
Restart Settlement
\end{minipage} \\
\midrule\noalign{}
\endhead
\bottomrule\noalign{}
\endlastfoot
Simple API call & 3 & Step 2 & \textbackslash{} \textasciitilde65s
(heartbeat) & \textasciitilde4h (manual) & 66\% (2/3 steps) & 0\% &
\textbackslash{} \textasciitilde90s total & Full re-execution \\
DeFi rebalance & 5 & Step 3 & \textasciitilde65s & \textbackslash{}
\textasciitilde4h & 60\% (3/5 steps) & 0\% & \textasciitilde85s total &
Full re-execution \\
Data pipeline & 10 & Step 7 & \textbackslash{} \textasciitilde65s &
\textasciitilde4h & 70\% (7/10 steps) & 0\% & \textbackslash{}
\textasciitilde70s total & Full re-execution \\
Long analysis & 20 & Step 15 & \textasciitilde65s & \textbackslash{}
\textasciitilde4h & 75\% (15/20 steps) & 0\% & \textasciitilde65s total
& Full re-execution \\
Complex pipeline & 50 & Step 42 & \textbackslash{} \textasciitilde65s &
\textasciitilde4h & 84\% (42/50 steps) & 0\% & \textbackslash{}
\textasciitilde60s total & Full re-execution \\
\end{longtable}
}

\textbf{Key insight:} With restart, failure cost is proportional to
total task length (all work is lost). With CAIRN, failure cost is
proportional to \emph{remaining} work only (completed checkpoints are
preserved). For a 50-step task failing at step 42, CAIRN preserves 84\%
of completed work and only re-executes the remaining 16\%.

\textbf{Economic comparison} (0.01 ETH escrow, 0.5\% protocol fee):

{\def\LTcaptype{none} % do not increment counter
\begin{longtable}[]{@{}
  >{\raggedright\arraybackslash}p{(\linewidth - 12\tabcolsep) * \real{0.1429}}
  >{\raggedright\arraybackslash}p{(\linewidth - 12\tabcolsep) * \real{0.1429}}
  >{\raggedright\arraybackslash}p{(\linewidth - 12\tabcolsep) * \real{0.1429}}
  >{\raggedright\arraybackslash}p{(\linewidth - 12\tabcolsep) * \real{0.1429}}
  >{\raggedright\arraybackslash}p{(\linewidth - 12\tabcolsep) * \real{0.1429}}
  >{\raggedright\arraybackslash}p{(\linewidth - 12\tabcolsep) * \real{0.1429}}
  >{\raggedright\arraybackslash}p{(\linewidth - 12\tabcolsep) * \real{0.1429}}@{}}
\toprule\noalign{}
\begin{minipage}[b]{\linewidth}\raggedright
Task Profile
\end{minipage} & \begin{minipage}[b]{\linewidth}\raggedright
Steps
\end{minipage} & \begin{minipage}[b]{\linewidth}\raggedright
Failure at
\end{minipage} & \begin{minipage}[b]{\linewidth}\raggedright
Original Agent (CAIRN)
\end{minipage} & \begin{minipage}[b]{\linewidth}\raggedright
Original Agent (Restart)
\end{minipage} & \begin{minipage}[b]{\linewidth}\raggedright
Wasted Gas (CAIRN)
\end{minipage} & \begin{minipage}[b]{\linewidth}\raggedright
Wasted Gas (Restart)
\end{minipage} \\
\midrule\noalign{}
\endhead
\bottomrule\noalign{}
\endlastfoot
DeFi rebalance & 5 & Step 3 & 0.00597 ETH (60\%) & 0 ETH & Steps 1-3:
\$0 & Steps 1-3: \textasciitilde\$0.003 \\
Data pipeline & 10 & Step 7 & 0.00697 ETH (70\%) & 0 ETH & Steps 1-7:
\$0 & Steps 1-7:\textbackslash{} \textasciitilde\$0.006 \\
Complex pipeline & 50 & Step 42 & 0.00836 ETH (84\%) & 0 ETH & Steps
1-42: \$0 & Steps 1-42: \textasciitilde\$0.035 \\
\end{longtable}
}

With CAIRN, the original agent is compensated for verified work. Without
CAIRN, 100\% of work and payment is lost on any failure.

\subsection{2.6 What CAIRN Is NOT}\label{what-cairn-is-not}

\begin{itemize}
\tightlist
\item
  \textbf{Not a new agent framework.} CAIRN wraps any existing framework
  --- LangGraph, Olas SDK, AgentKit, CrewAI, custom builds.
\item
  \textbf{Not a replacement for A2A or MCP.} Google's A2A protocol
  {[}12{]} handles agent discovery and communication. Anthropic's MCP
  {[}13{]} connects agents to tools. CAIRN handles what happens when
  those agents fail mid-task: detection, recovery, and settlement. These
  are complementary layers.
\item
  \textbf{Not a replacement for ERC-8183.} Virtuals' Agent Commerce
  Protocol implements ERC-8183 for the job lifecycle happy path (job
  creation → completion → payment). CAIRN handles the unhappy path
  (failure → classification → recovery → settlement). They compose.
\item
  \textbf{Not a centralized service.} Every state transition is enforced
  by the CAIRN state machine contract. No server. No admin key. No
  human-in-the-loop after task submission (the operator who submits the
  task may be a human, an agent, a DAO, or a contract --- see Section
  2.1).
\item
  \textbf{Not optional infrastructure.} The escrow condition makes
  record-writing mandatory --- agents cannot receive payment without
  completing the protocol.
\end{itemize}

\begin{center}\rule{0.5\linewidth}{0.5pt}\end{center}

\section{3. Design Philosophy}\label{design-philosophy}

\subsection{3.1 Classify by Recoverability, Not
Symptom}\label{classify-by-recoverability-not-symptom}

Prior research identifies 14+ failure modes in multi-agent systems
{[}1{]}, but existing taxonomies describe surface symptoms (``step
repetition,'' ``wrong tool selected'') without prescribing what to do
next. CAIRN's classification directly determines protocol behavior:

\textbf{LIVENESS failures} (\emph{F} = 0.70) --- the agent stopped
responding. A heartbeat was missed, a process crashed, or a network
partition occurred. These are almost always recoverable\textbackslash{}
(\textasciitilde92\% base rate): the task is not impossible, the agent
simply died. A fallback can resume from the last checkpoint immediately.

\textbf{RESOURCE failures} (\emph{F} = 0.30) --- the agent exhausted a
resource. Budget exceeded, deadline hit, API rate-limited, or context
window overflowed. These are partially recoverable (\textasciitilde48\%
base rate): success depends on whether sufficient budget and deadline
remain for the fallback.

\textbf{LOGIC failures} (\emph{F} = 0.00) --- the agent reasoned
incorrectly. Step repetition loops, hallucinated outputs, or
specification mismatches. These are rarely recoverable\textbackslash{}
(\textasciitilde8\% base rate): a fallback with the same task
specification will likely fail the same way. Setting \emph{F} = 0.00
routes all LOGIC failures directly to dispute.

This mapping is analogous to the foundational distinction between
\textbf{crash faults} and \textbf{Byzantine faults} in distributed
systems {[}8{]}. A crashed agent needs a different recovery path than an
agent producing wrong outputs. CAIRN operationalizes this insight for
the AI agent domain.

\textbf{Sub-class modulation.} Within each class, the failure type
provides additional signal that the fallback agent uses for off-chain
strategy selection (e.g., retry with different API key vs.~reduced
context vs.~alternative model), though it does not affect the on-chain
recovery score. The separation is deliberate: the score determines
\emph{whether} to attempt recovery (an on-chain decision requiring
determinism), while the failure type informs \emph{how} to attempt
recovery (an off-chain decision that benefits from richness).

{\def\LTcaptype{none} % do not increment counter
\begin{longtable}[]{@{}
  >{\raggedright\arraybackslash}p{(\linewidth - 6\tabcolsep) * \real{0.2500}}
  >{\raggedright\arraybackslash}p{(\linewidth - 6\tabcolsep) * \real{0.2500}}
  >{\raggedright\arraybackslash}p{(\linewidth - 6\tabcolsep) * \real{0.2500}}
  >{\raggedright\arraybackslash}p{(\linewidth - 6\tabcolsep) * \real{0.2500}}@{}}
\toprule\noalign{}
\begin{minipage}[b]{\linewidth}\raggedright
Failure Class
\end{minipage} & \begin{minipage}[b]{\linewidth}\raggedright
Failure Types
\end{minipage} & \begin{minipage}[b]{\linewidth}\raggedright
On-Chain Score Impact
\end{minipage} & \begin{minipage}[b]{\linewidth}\raggedright
Off-Chain Strategy Impact
\end{minipage} \\
\midrule\noalign{}
\endhead
\bottomrule\noalign{}
\endlastfoot
LIVENESS & HEARTBEAT\_MISS, PROCESS\_CRASH, NETWORK\_PARTITION & Same
(\emph{F} = 0.70) & Fallback uses same approach vs.~different node \\
RESOURCE & BUDGET\_EXCEEDED, DEADLINE\_HIT, RATE\_LIMIT,
CONTEXT\_OVERFLOW & Same (\emph{F} = 0.30) & Fallback uses different API
key vs.~smaller context vs.~reduced scope \\
LOGIC & HALLUCINATION, SPEC\_MISMATCH, STEP\_LOOP, WRONG\_TOOL & Same
(\emph{F} = 0.00) & N/A (routes to dispute) \\
\end{longtable}
}

Future protocol versions may introduce sub-class weights if production
data reveals that within-class recovery rate variance exceeds
between-class variance for specific failure types.

\subsection{3.2 Resume, Not Restart}\label{resume-not-restart}

Without checkpoints, a fallback agent must restart the entire task ---
wasting the original agent's completed work and the budget spent on it.
Checkpoints commit verified work after each subtask. On recovery, the
fallback reads the checkpoint list and resumes from the last verified
output. No restart from zero.

The theoretical foundation for this approach is the Chandy-Lamport
algorithm for distributed snapshots {[}8{]}, which proves that
consistent global state can be reconstructed from local checkpoints in a
distributed system. CAIRN adapts this: agents independently checkpoint
after each subtask (independent timing), but checkpoints are
schema-validated and IPFS-stored (coordinated verification) --- a
quasi-synchronous model.

\subsection{3.3 Permissionless
Enforcement}\label{permissionless-enforcement}

All enforcement functions (\texttt{checkLiveness}, \texttt{checkBudget},
\texttt{checkDeadline}) are public. Anyone can call them. No trusted
keeper is required.

The enforcement function only succeeds if the condition is actually
violated. False calls revert with no state change and no gas refund for
the caller. This makes the protocol permissionless --- any participant
can enforce liveness on any task, removing the dependency on centralized
keeper networks.

\subsection{3.4 Escrow as Forcing
Function}\label{escrow-as-forcing-function}

The escrow condition bootstraps participation without relying on
altruism. Agents cannot receive payment without completing the protocol
--- including writing the execution record. This creates a compounding
network effect:

\begin{verbatim}
More agents writing records → Richer intelligence layer
  → More accurate fallback selection → Higher recovery success rate
    → More agents integrating CAIRN → More agents writing records
\end{verbatim}

The loop starts from day one because the economic incentive is
immediate.

\subsection{3.5 Why On-Chain}\label{why-on-chain}

A legitimate question: why does agent failure recovery require a
blockchain? The answer is structural, not ideological.

\textbf{Escrow requires trustless settlement.} CAIRN distributes escrow
proportionally to verified checkpoints from two agents who do not trust
each other (the primary and fallback). A centralized coordinator could
modify checkpoint counts or settlement calculations. On-chain settlement
removes this trust dependency --- the escrow split is computed by an
immutable function visible to all parties.

\textbf{Permissionless enforcement requires public state.} Any
participant can call \texttt{checkLiveness} to trigger failure
detection. This is only possible when the enforcement function, the last
heartbeat timestamp, and the heartbeat interval are all publicly
readable on-chain. A centralized system could restrict who is allowed to
report failures.

\textbf{Record immutability requires append-only storage.} The execution
intelligence layer's value depends on records being tamper-proof. If a
centralized operator could modify failure records (e.g., deleting
records that make a preferred agent look bad), the intelligence layer
loses integrity. On-chain events are immutable.

\textbf{What does NOT need to be on-chain:} - Checkpoint data content →
IPFS (content-addressed, verifiable off-chain) - Intelligence queries →
The Graph / knowledge graph (off-chain indexing) - Pattern detection →
Off-chain analytics pipeline - Agent framework integration → SDK wraps
any framework locally

CAIRN places the minimum viable state on-chain (task state, escrow,
checkpoint counts, heartbeat timestamps) and keeps everything else
off-chain. This minimizes gas costs while preserving the trust
properties that motivate blockchain use.

\begin{center}\rule{0.5\linewidth}{0.5pt}\end{center}

\section{4. Key Design Decisions}\label{key-design-decisions}

\subsection{4.1 Checkpoint Protocol}\label{checkpoint-protocol}

\textbf{Write flow:}

\begin{verbatim}
Agent completes subtask N
→ Agent writes output to IPFS → receives content-addressed CID
→ Agent calls commitCheckpointBatch(taskId, count, merkleRoot, latestCID)
→ CAIRN stores Merkle root and latest CID, increments checkpoint count
→ Event: CheckpointBatchCommitted(taskId, count, merkleRoot)
\end{verbatim}

\textbf{Recovery read flow:}

\begin{verbatim}
Fallback agent receives task state:
  - checkpoint CIDs: [CID_0, CID_1, CID_2]  (from events/IPFS)
  - next subtask index: 3                     (resume here)
  - remaining budget, remaining deadline
→ Fallback reads CID_2 from IPFS → subtask 2 output
→ Fallback begins subtask 3 using subtask 2 output as input
\end{verbatim}

Checkpoints use \textbf{Merkle tree batching} for gas efficiency:
multiple checkpoint CIDs are batched off-chain, and only the Merkle root
is committed on-chain. Individual checkpoints can be verified via Merkle
proof when needed (e.g., during dispute). This reduces gas costs by
approximately 95\% for tasks with many checkpoints compared to per-CID
storage {[}see Section 6.5{]}.

\textbf{Incentive alignment:} Agents are paid proportionally to verified
checkpoint count. More checkpoints written means more partial payment on
failure. This incentivizes frequent, honest checkpointing --- the
original agent has a direct financial interest in making their work
resumable.

\subsection{4.1.1 Checkpoint Portability}\label{checkpoint-portability}

\begin{quote}
\textbf{Scope.} CAIRN's recovery guarantees apply \emph{fully} to
structured-pipeline tasks (data fetches, API calls, multi-step
computations, stateful queries --- see ``Fully portable'' and ``Portable
with context'' rows below) and \emph{partially} to reasoning-heavy tasks
(chain-of-thought, planning with backtracking ---
``Framework-dependent''). For the framework-dependent class, only
output-level checkpoints are portable; the fallback resumes the next
subtask from the output of the previous one, losing any implicit
reasoning state. The 23.46\% misrouting result and the §6.6 economic
analysis are calibrated for the first two classes, which we estimate
cover \textasciitilde80-90\% of current on-chain agent workloads based
on the task-type distribution observed in Olas and Virtuals deployments.
Reasoning-heavy workloads receive degraded recovery rather than no
recovery; quantifying the exact degradation requires empirical study,
which is reserved for v3 (Section 10.1, open research questions).
\end{quote}

A checkpoint's value depends on whether a different agent ---
potentially running a different framework --- can meaningfully resume
from it. We define portability formally:

\textbf{Definition (Semantic Portability).} A checkpoint\textbackslash{}
\emph{C}\textsubscript{i} for subtask \emph{i} is \emph{semantically
portable} from framework\textbackslash{} \emph{F}\textsubscript{1} to
framework\textbackslash{} \emph{F}\textsubscript{2} if and only if an
agent running\textbackslash{} \emph{F}\textsubscript{2} can produce a
correct output for subtask \emph{i+1} given only\textbackslash{}
\emph{C}\textsubscript{i} and the task specification \emph{S}, without
access to\textbackslash{} \emph{F}\textsubscript{1}'s internal state.

Portability depends on the checkpoint's \textbf{context completeness}
--- whether the checkpoint contains all information needed to continue
the task:

{\def\LTcaptype{none} % do not increment counter
\begin{longtable}[]{@{}
  >{\raggedright\arraybackslash}p{(\linewidth - 6\tabcolsep) * \real{0.2500}}
  >{\raggedright\arraybackslash}p{(\linewidth - 6\tabcolsep) * \real{0.2500}}
  >{\raggedright\arraybackslash}p{(\linewidth - 6\tabcolsep) * \real{0.2500}}
  >{\raggedright\arraybackslash}p{(\linewidth - 6\tabcolsep) * \real{0.2500}}@{}}
\toprule\noalign{}
\begin{minipage}[b]{\linewidth}\raggedright
Portability Class
\end{minipage} & \begin{minipage}[b]{\linewidth}\raggedright
Description
\end{minipage} & \begin{minipage}[b]{\linewidth}\raggedright
Context Requirement
\end{minipage} & \begin{minipage}[b]{\linewidth}\raggedright
Examples
\end{minipage} \\
\midrule\noalign{}
\endhead
\bottomrule\noalign{}
\endlastfoot
\textbf{Fully portable} & Checkpoint output is self-contained & Output
is the complete result of subtask \emph{i}; no implicit state needed &
Data fetches, API calls, file transformations, on-chain reads \\
\textbf{Portable with context} & Checkpoint output is complete when
combined with explicit metadata & Output plus metadata fields (e.g.,
accumulated state, configuration) & Multi-step computations, stateful
API sessions, paginated queries \\
\textbf{Framework-dependent} & Checkpoint requires internal framework
state to interpret & Implicit reasoning chain, model context window,
conversation history & Chain-of-thought reasoning, multi-turn dialogue,
planning with backtracking \\
\end{longtable}
}

CAIRN's checkpoint schema targets the first two classes by requiring
explicit context serialization:

\begin{Shaded}
\begin{Highlighting}[]
\FunctionTok{\{}
  \DataTypeTok{"task\_id"}\FunctionTok{:} \StringTok{"0x..."}\FunctionTok{,}
  \DataTypeTok{"subtask\_index"}\FunctionTok{:} \DecValTok{3}\FunctionTok{,}
  \DataTypeTok{"output"}\FunctionTok{:} \FunctionTok{\{} \ErrorTok{...} \FunctionTok{\},}
  \DataTypeTok{"context"}\FunctionTok{:} \FunctionTok{\{}
    \DataTypeTok{"accumulated\_state"}\FunctionTok{:} \FunctionTok{\{} \ErrorTok{...} \FunctionTok{\},}
    \DataTypeTok{"dependencies"}\FunctionTok{:} \OtherTok{[}\StringTok{"subtask\_0\_cid"}\OtherTok{,} \StringTok{"subtask\_1\_cid"}\OtherTok{,} \StringTok{"subtask\_2\_cid"}\OtherTok{]}\FunctionTok{,}
    \DataTypeTok{"metadata"}\FunctionTok{:} \FunctionTok{\{} \ErrorTok{...} \FunctionTok{\}}
  \FunctionTok{\},}
  \DataTypeTok{"schema\_version"}\FunctionTok{:} \StringTok{"1.0"}
\FunctionTok{\}}
\end{Highlighting}
\end{Shaded}

The \texttt{context} field is the portability mechanism: any agent can
reconstruct the necessary state by reading the output and context
fields, without needing the original agent's internal representation.
For framework-dependent tasks (reasoning chains with implicit state),
operators should decompose the task into subtasks whose outputs are
self-contained --- effectively converting framework-dependent
checkpoints into the ``portable with context'' class.

\subsection{4.1.3 Data Availability}\label{data-availability}

\textbf{Checkpoint availability.} CAIRN requires checkpoint data
availability for two windows: (1) the task duration (for fallback
resumption) and (2) the dispute period (for arbiter evidence).
Availability is ensured via:

\begin{itemize}
\tightlist
\item
  \textbf{Protocol-level pinning:} The CAIRN SDK pins all checkpoint
  CIDs to a configurable pinning service (default: Pinata) on commit.
  The protocol fee (Section 6.2) covers pinning costs for the task
  duration plus dispute period.
\item
  \textbf{Operator-level redundancy:} Operators may specify additional
  pinning services at task submission.
\item
  \textbf{Availability fallback:} If a checkpoint CID is unretrievable
  during recovery, the fallback resumes from the last available
  checkpoint (potentially losing work between the last available and the
  actual last checkpoint). If no checkpoints are retrievable, the task
  routes to DISPUTED.
\end{itemize}

Future versions may integrate Filecoin storage deals for cryptographic
availability guarantees, or EIP-4844 blob storage for short-lived
checkpoint data.

\subsection{4.2 Task Type Taxonomy}\label{task-type-taxonomy}

Every routing decision depends on \texttt{task\_type}. The taxonomy is
hierarchical:

\begin{verbatim}
task_type = domain.operation
\end{verbatim}

Examples: \texttt{defi.price\_fetch}, \texttt{defi.trade\_execute},
\texttt{data.report\_generate}, \texttt{governance.vote\_delegate},
\texttt{compute.model\_inference}.

Agents declare supported task types in their ERC-8004 identity record.
Fallback matching follows a precedence order: (1) exact match on
\texttt{domain.operation} with highest reputation and available stake;
(2) domain-level match with highest reputation; (3) no match → DISPUTED
immediately.

\subsection{4.3 Adaptive Liveness
Interval}\label{adaptive-liveness-interval}

A fixed heartbeat interval is incorrect for all task types. A 30-second
API call and a 3-hour analysis task require different liveness
requirements.

\begin{verbatim}
heartbeat_interval = operator_declared_value
subject to:
  min = 30 seconds (Base block time ≈ 2s → 15 blocks)
  max = task_deadline / 4
  default = min(task_deadline / 10, 300 seconds)
\end{verbatim}

This ensures at least 10 liveness signals per task by default, with a
5-minute cap per interval. The v1 \texttt{isStale(taskId)} check
triggers failure detection only after \textbf{two} consecutive missed
intervals (i.e.,
\texttt{block.timestamp\ \textgreater{}\ lastHeartbeat\ +\ 2\ ×\ heartbeatInterval}),
which widens the effective liveness window to absorb transient RPC or
sequencer hiccups without false-positive failure detection. A single
late heartbeat does not fail the task; two consecutive misses do.

\textbf{Progress detection.} Heartbeats confirm the agent process is
alive but not that it is making progress. A stuck agent (infinite loop,
hung API call) continues to heartbeat while producing no checkpoints.
CAIRN detects this via a progress timeout: if no new checkpoint is
committed within
\texttt{max(heartbeat\_interval\ ×\ 4,\ expected\_subtask\_duration)},
the protocol considers the agent stalled. The operator can configure
\texttt{expected\_subtask\_duration} at task submission (default:
\texttt{deadline\ /\ total\_expected\_subtasks}).

Stalled agents are treated as RESOURCE failures (the agent has consumed
time without producing output). The progress timeout is enforced by a
public \texttt{checkProgress(taskId)} function analogous to
\texttt{checkLiveness}.

\textbf{Gameability bound.} An adversarial worker emitting heartbeats
while producing no useful work is bounded by the \emph{tighter} of the
two timeouts. Specifically: between any two consecutive checkpoints, an
agent can stall for at most
\texttt{max(4\ ×\ heartbeat\_interval,\ expected\_subtask\_duration)}
before \texttt{checkProgress} fires. With the default
\texttt{heartbeat\_interval\ =\ min(deadline\ /\ 10,\ 300s)} and the
default
\texttt{expected\_subtask\_duration\ =\ deadline\ /\ n\_subtasks}, the
worst-case stall window per subtask is the larger of: -
\texttt{4\ ×\ deadline\ /\ 10\ =\ 0.4\ ×\ deadline} (heartbeat-bound) -
\texttt{deadline\ /\ n\_subtasks} (subtask-bound)

For a typical task with \texttt{n\_subtasks\ ≥\ 5}, the subtask bound is
tighter, capping stall time at ≤ 20\% of deadline per subtask. For a
task with very few subtasks (\texttt{n\ ≤\ 3}), the heartbeat bound
dominates and an adversarial worker can extract up to\textbackslash{}
\textasciitilde24\% of deadline cost-free (slightly under the
\texttt{4×\ heartbeat\_interval} threshold). Operators submitting
low-subtask tasks should set \texttt{expected\_subtask\_duration}
explicitly rather than relying on the default; setting it to
\texttt{1.5\ ×\ heartbeat\_interval} reduces the gameable window to
\textasciitilde6\% of deadline. This trade-off --- tighter progress
windows reduce gameability but increase false-positive stall detections
on legitimately variable subtasks --- is left to operator configuration
rather than fixed by the protocol.

\subsection{4.4 Fallback Pool Admission
Control}\label{fallback-pool-admission-control}

Open registration creates a vulnerability: malicious or unreliable
agents could register broadly, accept recovery assignments, and collect
partial payment without completing work.

\textbf{Two-gate admission:}

\textbf{Gate 1 --- Reputation:} Minimum ERC-8004 reputation score for
the declared task type. Default threshold: score ≥ 50 on a 0-100 scale.

\textbf{Gate 2 --- Stake:} Deposit proportional to maximum eligible
escrow. Default: \texttt{min\_stake\ =\ max\_eligible\_escrow\ ×\ 0.1}.
If the fallback agent fails without completing any checkpoints, the full
stake is slashed and distributed to the operator.

\textbf{Optional: Olas Mech Marketplace integration.} When no internal
fallback is available, CAIRN queries the Olas Mech Marketplace {[}5{]}
for eligible agents by task capability, filtered by minimum reputation
(85\% success rate). This could extend the fallback pool to
Olas's\textbackslash{} \textasciitilde2,000 deployed mechs (≈500 active
daily), subject to integration of CAIRN's checkpoint schema with the
Olas execution model.

\subsection{4.5 Arbiter Design}\label{arbiter-design}

The arbiter role is itself an agent service. Arbiter agents register in
CAIRN with a stake proportional to the maximum dispute value they can
rule on (\texttt{min\_arbiter\_stake\ =\ max\_dispute\_value\ ×\ 0.2}).
They read public execution records, submit rulings, and earn fees (3\%
of dispute value).

Sybil resistance is economic: incorrect rulings (detectable via on-chain
evidence) result in stake slashing. The stake at risk (20\%) exceeds the
fee earned (3\%), making collusion uneconomical \emph{for disputes where
detectability holds} --- see Section 7.5 Proposition 3 for the formal
analysis and the explicit acknowledgment that detection probability is
high for LIVENESS and RESOURCE disputes (on-chain evidence) but
materially lower for LOGIC disputes (where ``incorrect'' is a reasoning
judgment). The v2 specification ships single-tier arbitration with this
caveat; an appeals layer (in the style of Kleros {[}19{]}) is reserved
for v3.

A commit-reveal scheme prevents front-running of dispute rulings. If no
arbiter rules within the timeout (default 7 days), escrow auto-refunds
to the operator.

\begin{center}\rule{0.5\linewidth}{0.5pt}\end{center}

\section{5. Execution Intelligence
Layer}\label{execution-intelligence-layer}

\subsection{5.1 What Gets Written}\label{what-gets-written}

Record-writing is mandatory for escrow settlement --- not optional.

\textbf{On FAILED --- Failure Record (IPFS):}

\begin{Shaded}
\begin{Highlighting}[]
\FunctionTok{\{}
  \DataTypeTok{"record\_type"}\FunctionTok{:} \StringTok{"failure"}\FunctionTok{,}
  \DataTypeTok{"task\_id"}\FunctionTok{:} \StringTok{"0x..."}\FunctionTok{,}
  \DataTypeTok{"agent\_id"}\FunctionTok{:} \StringTok{"erc8004://base/0x..."}\FunctionTok{,}
  \DataTypeTok{"task\_type"}\FunctionTok{:} \StringTok{"defi.price\_fetch"}\FunctionTok{,}
  \DataTypeTok{"failure\_class"}\FunctionTok{:} \StringTok{"LIVENESS"}\FunctionTok{,}
  \DataTypeTok{"failure\_type"}\FunctionTok{:} \StringTok{"HEARTBEAT\_MISS"}\FunctionTok{,}
  \DataTypeTok{"checkpoint\_count\_at\_failure"}\FunctionTok{:} \DecValTok{3}\FunctionTok{,}
  \DataTypeTok{"cost\_at\_failure"}\FunctionTok{:} \StringTok{"0.0023 ETH"}\FunctionTok{,}
  \DataTypeTok{"budget\_remaining\_pct"}\FunctionTok{:} \FloatTok{0.42}\FunctionTok{,}
  \DataTypeTok{"deadline\_remaining\_pct"}\FunctionTok{:} \FloatTok{0.31}\FunctionTok{,}
  \DataTypeTok{"recovery\_score"}\FunctionTok{:} \FloatTok{0.47}\FunctionTok{,}
  \DataTypeTok{"timestamp"}\FunctionTok{:} \DecValTok{1742000000}
\FunctionTok{\}}
\end{Highlighting}
\end{Shaded}

\textbf{On RESOLVED --- Resolution Record (IPFS):}

\begin{Shaded}
\begin{Highlighting}[]
\FunctionTok{\{}
  \DataTypeTok{"record\_type"}\FunctionTok{:} \StringTok{"resolution"}\FunctionTok{,}
  \DataTypeTok{"task\_id"}\FunctionTok{:} \StringTok{"0x..."}\FunctionTok{,}
  \DataTypeTok{"states\_traversed"}\FunctionTok{:} \OtherTok{[}\StringTok{"RUNNING"}\OtherTok{,} \StringTok{"FAILED"}\OtherTok{,} \StringTok{"RECOVERING"}\OtherTok{,} \StringTok{"RESOLVED"}\OtherTok{]}\FunctionTok{,}
  \DataTypeTok{"task\_type"}\FunctionTok{:} \StringTok{"defi.price\_fetch"}\FunctionTok{,}
  \DataTypeTok{"total\_cost"}\FunctionTok{:} \StringTok{"0.0041 ETH"}\FunctionTok{,}
  \DataTypeTok{"original\_checkpoint\_count"}\FunctionTok{:} \DecValTok{3}\FunctionTok{,}
  \DataTypeTok{"fallback\_checkpoint\_count"}\FunctionTok{:} \DecValTok{2}\FunctionTok{,}
  \DataTypeTok{"escrow\_split"}\FunctionTok{:} \FunctionTok{\{}
    \DataTypeTok{"original\_agent"}\FunctionTok{:} \StringTok{"0.006 ETH"}\FunctionTok{,}
    \DataTypeTok{"fallback\_agent"}\FunctionTok{:} \StringTok{"0.004 ETH"}\FunctionTok{,}
    \DataTypeTok{"protocol\_fee"}\FunctionTok{:} \StringTok{"0.00005 ETH"}
  \FunctionTok{\}}
\FunctionTok{\}}
\end{Highlighting}
\end{Shaded}

\subsection{5.2 What Agents Query}\label{what-agents-query}

\textbf{Pre-task intelligence (Phase 1 --- Initialization):}

{\def\LTcaptype{none} % do not increment counter
\begin{longtable}[]{@{}
  >{\raggedright\arraybackslash}p{(\linewidth - 4\tabcolsep) * \real{0.2800}}
  >{\raggedright\arraybackslash}p{(\linewidth - 4\tabcolsep) * \real{0.3600}}
  >{\raggedright\arraybackslash}p{(\linewidth - 4\tabcolsep) * \real{0.3600}}@{}}
\toprule\noalign{}
\begin{minipage}[b]{\linewidth}\raggedright
Query
\end{minipage} & \begin{minipage}[b]{\linewidth}\raggedright
Returns
\end{minipage} & \begin{minipage}[b]{\linewidth}\raggedright
Purpose
\end{minipage} \\
\midrule\noalign{}
\endhead
\bottomrule\noalign{}
\endlastfoot
Failure patterns by task type & Ranked failure types with frequency &
Operator assesses risk before committing escrow \\
Cost distribution & P25/P50/P75/P95 total cost & Operator sets realistic
budget cap \\
Agent success rates & Agents ranked by success rate for task type &
Select best-fit primary agent \\
Known-bad conditions & Time windows or APIs correlated with failures &
Avoid scheduling during high-risk windows \\
\end{longtable}
}

\textbf{Fallback selection intelligence (on FAILED → RECOVERING):}

\begin{verbatim}
Input:  task_type, remaining_budget, remaining_deadline
Query:  Eligible agents sorted by:
        1. Success rate on this task_type
        2. ERC-8004 reputation score
        3. Stake deposited
        4. Current availability
Filter: Admission threshold (min reputation + active stake)
Output: Ranked list of eligible fallback agents
\end{verbatim}

\subsection{5.3 Network Effects}\label{network-effects}

The execution history is on-chain and publicly indexable, but
replicating its accumulation requires matching CAIRN's escrow-mandated
write volume --- a competing protocol would need to independently
generate comparable failure data, which requires comparable task
throughput. A competitor can copy the protocol code. They cannot copy
the accumulated records --- the failure patterns, the agent performance
data, the cost distributions --- without first reproducing the economic
activity that produced them. This creates a defensible network effect
where each new agent failure makes the protocol more valuable for every
future agent.

\textbf{Quantitative model.} The intelligence layer's utility for a
given task type \emph{τ} is a function of the number of recorded failure
and resolution events \emph{n}\textsubscript{τ}:

\begin{verbatim}
Pattern confidence:     P(τ, n) = 1 −\\ e^{−n/k}
Fallback accuracy:      F(τ, n) = F_0 + (F_max − F_0) × (1 − e^{−n/m})
\end{verbatim}

Where: - \emph{k} = minimum records for 63\% pattern confidence
(estimated: \emph{k} ≈ 30 per task type). The value \emph{k} ≈ 30
follows from the sample size formula for a one-proportion \emph{z}-test:
\emph{n} =\textbackslash{} (\emph{z}\textsubscript{α/2})² ×
\emph{p}(1−\emph{p}) / \emph{d}², where \emph{p} = 0.5 (failure rate),
\emph{d} = 0.10 (precision), \emph{α} = 0.05 → \emph{n} = 1.96² × 0.25 /
0.01 ≈ 96 total records. With a 45\% LIVENESS proportion, we
need\textbackslash{} \textasciitilde96 × 0.45 ≈ 43 LIVENESS records ---
or roughly 30 records per class when amortized across the three classes.
- \emph{F}\textsubscript{0} = baseline fallback success rate without
intelligence (random selection from pool) -\textbackslash{}
\emph{F}\textsubscript{max} = maximum fallback success rate with full
intelligence - \emph{m} = records needed for 63\% of maximum improvement
(estimated: \emph{m} ≈ 100 per task type)

\textbf{Minimum viable intelligence thresholds:}

{\def\LTcaptype{none} % do not increment counter
\begin{longtable}[]{@{}
  >{\raggedright\arraybackslash}p{(\linewidth - 4\tabcolsep) * \real{0.3333}}
  >{\raggedright\arraybackslash}p{(\linewidth - 4\tabcolsep) * \real{0.3333}}
  >{\raggedright\arraybackslash}p{(\linewidth - 4\tabcolsep) * \real{0.3333}}@{}}
\toprule\noalign{}
\begin{minipage}[b]{\linewidth}\raggedright
Records per Task Type
\end{minipage} & \begin{minipage}[b]{\linewidth}\raggedright
Pattern Confidence
\end{minipage} & \begin{minipage}[b]{\linewidth}\raggedright
Practical Utility
\end{minipage} \\
\midrule\noalign{}
\endhead
\bottomrule\noalign{}
\endlastfoot
\emph{n} = 10 & 28\% & Preliminary signals; not actionable \\
\emph{n} = 30 & 63\% & Dominant failure type identifiable \\
\emph{n} = 100 & 96\% & Reliable failure pattern distribution; fallback
ranking meaningful \\
\emph{n} = 300 & \textgreater99.99\% & Statistical significance for
time-based and agent-specific patterns \\
\end{longtable}
}

These estimates follow from the exponential model: at \emph{n} =
\emph{k}, confidence is 1 −\textbackslash{} \emph{e}\textsuperscript{−1}
≈ 0.63. The model predicts diminishing returns --- the first 100 records
per task type provide the majority of intelligence value, making the
cold-start problem bounded rather than open-ended.

\textbf{Cold-start bootstrap.} The initial 6 task types are
\texttt{defi.price\_fetch}, \texttt{defi.trade\_execute},
\texttt{data.report\_generate}, \texttt{governance.vote\_delegate},
\texttt{compute.model\_inference}, and a reserved \texttt{generic.*}
catch-all. Using the simulation's 50\% failure rate (base literature
value {[}3{]}, modulated downward to 36.8\% observed in Section 10.1
once complexity and skill factors apply): at a protocol-wide throughput
of \textbf{10 tasks/day}, 100 tasks per task type (≈ 60 days of
operation) produces\textbackslash{} \textasciitilde50 failure records
per type --- sufficient for 81\% pattern confidence (1 −\textbackslash{}
e\textsuperscript{−50/30}). Reaching \emph{k} = 30 records per type
(63\% confidence) takes \textasciitilde36 days at this rate. When
decomposed per failure class, the LOGIC path saturates
last\textbackslash{} (\textasciitilde180 days to \emph{k} = 30 LOGIC
records per type), but this does not block protocol utility: LOGIC's
recovery score is 0 by construction (Section 6.4), so it routes to
dispute without needing pattern confidence for the recover/dispute
decision. For the per-type aggregate decision (recover vs.~dispute),
minimum viable intelligence is reached in \textasciitilde60 days.

\begin{center}\rule{0.5\linewidth}{0.5pt}\end{center}

\section{6. Economic Model}\label{economic-model}

\subsection{6.1 Escrow Split Rule}\label{escrow-split-rule}

On RESOLVED, escrow is distributed proportionally to verified work:

\begin{verbatim}
protocol_fee      = escrow × fee_bps / 10000
distributable     = escrow - protocol_fee
original_share    = distributable × (original_checkpoints / total_checkpoints)
fallback_share    = distributable × (fallback_checkpoints / total_checkpoints)
\end{verbatim}

If no recovery occurred (original agent completed solo): 100\% of
distributable to original agent.

\subsection{6.2 Protocol Fee}\label{protocol-fee}

\begin{itemize}
\tightlist
\item
  Default: 50 basis points (0.5\%) of escrow on settlement
\item
  Collected on every RESOLVED state transition
\item
  Configurable in the v2 upgradeable implementation (range: 0-500 bps);
  the v1 testnet contract declares \texttt{protocolFeeBps} as
  \texttt{constant}, so fee changes on v1 require a code redeployment
  rather than a parameter set.
\end{itemize}

\subsection{6.3 Stake Requirements}\label{stake-requirements}

{\def\LTcaptype{none} % do not increment counter
\begin{longtable}[]{@{}
  >{\raggedright\arraybackslash}p{(\linewidth - 6\tabcolsep) * \real{0.1154}}
  >{\raggedright\arraybackslash}p{(\linewidth - 6\tabcolsep) * \real{0.3077}}
  >{\raggedright\arraybackslash}p{(\linewidth - 6\tabcolsep) * \real{0.3269}}
  >{\raggedright\arraybackslash}p{(\linewidth - 6\tabcolsep) * \real{0.2500}}@{}}
\toprule\noalign{}
\begin{minipage}[b]{\linewidth}\raggedright
Role
\end{minipage} & \begin{minipage}[b]{\linewidth}\raggedright
Min Stake (v2)
\end{minipage} & \begin{minipage}[b]{\linewidth}\raggedright
Slash Condition
\end{minipage} & \begin{minipage}[b]{\linewidth}\raggedright
Slash Amount
\end{minipage} \\
\midrule\noalign{}
\endhead
\bottomrule\noalign{}
\endlastfoot
Fallback agent & 10\% of max eligible escrow & Fails without completing
any checkpoints & 100\% of stake \\
Arbiter & 20\% of max ruleable dispute & Incorrect ruling (detectable
via evidence) & 50\% of stake \\
\end{longtable}
}

\textbf{Derivation of the 10\% fallback stake.} The minimum stake must
satisfy two simultaneous constraints. First, it must make
zero-checkpoint failure \emph{economically irrational} for any agent
whose ex-ante success rate exceeds the admission floor --- i.e., the
expected loss from staking exceeds the expected gain from an unjustified
acceptance. Combining the fallback acceptance condition (Proposition 2
in §7.5) with the slash mechanism: a fallback with success rate \emph{s}
faces expected gain
\texttt{s\ ·\ E\_r\ ·\ (c\_f/c\_total)\ ·\ (1−f)\ −\ (1−s)\ ·\ p\_0\ ·\ stake}.
For the admission floor \emph{s} = 0.5 (reputation gate),
zero-checkpoint failure probability\textbackslash{}
\emph{p}\textsubscript{0} ≈ 0.3 (empirical estimate from the simulation
ground truth where failures cluster early), and the typical case
\texttt{(c\_f/c\_total)\ ≈\ 0.4} (fallback takes\textbackslash{}
\textasciitilde40\% of work on recovery), the break-even stake is
approximately
\texttt{0.5\ ×\ 0.4\ ×\ E\_r\ /\ (0.5\ ×\ 0.3)\ ≈\ 1.33\ ×\ E\_r}. Since
\texttt{E\_r\ ≤\ E}, a stake of 10-15\% of \emph{maximum eligible
escrow} across the agent's task portfolio satisfies this break-even with
margin. Second, the stake must be small enough that capable agents will
actually post it --- 10\% is large enough to deter Sybils but small
enough that a reputable fallback can serve many tasks without locking
excessive capital. Empirically, 10\% is the floor used by similar
staked-fallback designs (e.g., Olas's mech registration); CAIRN inherits
this convention.

\textbf{Derivation of the 20\% arbiter stake.} The arbiter stake must
dominate the arbiter fee by a factor large enough to make collusion
uneconomical. With fee φ = 3\% of dispute value and slash amount 50\% of
stake on detection, the no-collusion condition is
\texttt{bribe\ \textless{}\ φ\ +\ p\_d\ ×\ (slash/2\ ×\ stake)\ =\ 0.03V\ +\ p\_d\ ×\ stake/2}.
Setting \texttt{stake\ =\ 0.20V} and \texttt{p\_d\ =\ 0.8}
(high-detection regime, on-chain evidence) yields
\texttt{bribe\ \textless{}\ 0.03V\ +\ 0.08V\ =\ 0.11V} --- the minimum
bribe must exceed 11\% of dispute value. For the LOGIC-dispute regime
where \emph{p}\textsubscript{d} ≈ 0.4, the same stake yields
\texttt{bribe\ \textless{}\ 0.05V}, which is the residual structural
risk discussed in §7.5 and §10.3. A higher stake (e.g., 30\%) would
tighten the LOGIC-dispute bound but at the cost of arbiter participation
--- 20\% is the compromise that keeps arbiter capital cost roughly equal
to\textbackslash{} \textasciitilde6-7× the per-dispute fee, which we
judge is the lowest acceptable risk-to-reward ratio that still attracts
professional arbiters. The v1 testnet uses 15\% pending the v2 upgrade.

The v1 testnet deployment currently enforces a \textbf{15\% arbiter
stake} (plus a 0.15 ETH absolute floor) rather than 20\%. The v2 upgrade
raises the ratio to 20\% to strengthen the incentive analysis in Section
7.5. Under v1 parameters, the Proposition 3 inequality becomes
\texttt{bribe\ \textless{}\ 0.03V\ +\ p\_d\ ×\ 0.075V}; the qualitative
conclusion (honest ruling rational under realistic detection
probabilities) is preserved, but the margin is tighter.

\subsection{\texorpdfstring{6.4 Recovery Score Formula \emph{(v2
specification)}}{6.4 Recovery Score Formula (v2 specification)}}\label{recovery-score-formula-v2-specification}

\textbf{v2 formula (simulation-validated, this paper's headline
result):}

\begin{verbatim}
r =\\ F^a × B^b ×\\ D^c
\end{verbatim}

Where \emph{r} is the recovery score, expanded with default parameters
as:

\begin{verbatim}
r = F^0.80 ×\\ B^0.35 × D^0.15
\end{verbatim}

Where: - \emph{F} = \texttt{failure\_class\_weight}: LIVENESS = 0.70,
RESOURCE = 0.30, LOGIC = 0.00 - \emph{B} =
\texttt{budget\_remaining\_pct}: (budget\_cap - cost\_accrued) /
budget\_cap - \emph{D} = \texttt{deadline\_remaining\_pct}: (deadline -
current\_block) / (deadline - start\_block) - (\emph{a}, \emph{b},
\emph{c}) = governance-adjustable exponents; default (0.80, 0.35, 0.15)

\textbf{Three-tier routing (v2):} - Score ≥ 0.40 → \textbf{RECOVERING}
(high confidence --- automatic fallback, full remaining budget) - 0.35 ≤
Score \textless{} 0.40 → \textbf{RECOVERING (reduced scope)} (attempt
with constraints --- fallback receives capped budget) - Score
\textless{} 0.35 → \textbf{DISPUTED} (requires arbiter resolution)

\begin{quote}
\textbf{v1 interim formula (current testnet deployment).} The v1
contract on Base Sepolia implements the pre-calibration linear formula
\texttt{r\ =\ 0.5·F\ +\ 0.3·B\ +\ 0.2·D} with class weights (0.90, 0.50,
0.10) and a single binary threshold at 0.30 --- the 47.56\%-misrouting
baseline called ``Eq1-current'' in Section 10.1. The calibration work in
this paper motivates the v2 upgrade path: replace the linear formula
with the multiplicative formula above, move class weights to (0.70,
0.30, 0.00), and introduce the three-tier threshold band (0.40 / 0.35).
The migration is governance-gated and does not require a state-breaking
upgrade, because the score formula is already isolated behind the
\texttt{IRecoveryRouter} interface in the v1 contract architecture.
\end{quote}

The three-tier model enables graduated recovery: high-confidence
failures get full resources, medium-confidence failures get a
constrained attempt before escalating to dispute, and low-confidence
failures go directly to arbitration.

\textbf{Why multiplicative.} The formula uses a product rather than a
weighted sum because recovery success depends on \emph{all} factors
being adequate simultaneously. If budget is zero, recovery is impossible
regardless of failure type or deadline. If the failure is a LOGIC error
(\emph{F} = 0.00), no amount of budget or time helps. The multiplicative
structure captures this ``any-factor-kills-it'' dynamic: when any input
approaches zero, the score approaches zero --- matching empirical
recovery dynamics.

This design choice is empirically validated. Monte Carlo simulation
across 100,000 synthetic task-failure events per run (seed=42,
reproducible via \texttt{python3\ -m\ simulation.run\_eq4})
systematically compared four formula structures across 16 experiments:
(1) linear weighted sum --- optimal 33.81\% misrouting (Run 1, 362 grid
points); (2) piecewise-linear with \emph{B}×\emph{D} interaction ---
33.17\% (Run 2); (3) 5-variable linear with complexity and skill inputs
--- 32.78\% (Run 3); and (4) multiplicative --- \textbf{23.46\%} (Run 4,
2,646 grid points). The first three formulas converge to a
\textasciitilde33\% misrouting floor --- a structural ceiling intrinsic
to additive formulas, confirmed across 3,008 configurations. The
multiplicative formula breaks through to 23.46\%, within \textbf{0.93
percentage points of the Bayes-optimal theoretical minimum (22.53\%)}
--- capturing 96\% of achievable improvement. A hybrid α-sweep (11
ratios from α=0.0 pure multiplicative to α=1.0 pure linear) confirms
that pure multiplicative is strictly optimal: misrouting increases
monotonically with α (23.46\% at α=0, 24.57\% at α=0.1, 26.27\% at
α=0.5, 35.07\% at α=1.0). Cross-task-type leave-one-out validation shows
23.39\% ± 0.36\% across five task types --- best generalization of any
formula tested. Most importantly, the confusion matrix pivots:
\textbf{FULL-tier false positives drop from 22.3\% (Eq1 linear) to 7.9\%
(Eq4 multiplicative) --- a 65\% reduction in wasted recovery attempts.}
Full methodology, per-experiment findings, and confusion matrices are
documented in Section 10.1; raw results in
\texttt{simulation/RESULTS\_EQ4.md}.

\textbf{Exponent rationale.} The failure class exponent \emph{a} = 0.80
makes \emph{F} the dominant factor: a LIVENESS failure (\emph{F} = 0.70)
produces\textbackslash{} \emph{F}\^{}0.80 ≈ 0.752, while a RESOURCE
failure (\emph{F} = 0.30) produces \emph{F}\^{}0.80 ≈ 0.382 --- roughly
a 2× separation (precise ratio 1.97). The sub-linear exponent provides
diminishing returns above \emph{F} = 0.5, preventing the class signal
from overwhelming resource signals. The budget exponent \emph{b} = 0.35
assigns moderate influence: 50\% budget remaining yields\textbackslash{}
\emph{B}\^{}0.35 = 0.79, while 10\% yields \emph{B}\^{}0.35 = 0.47 --- a
meaningful but not catastrophic penalty. The deadline exponent \emph{c}
= 0.15 assigns the least weight: in the multiplicative context, deadline
contributes through the product interaction (low deadline × low budget
is catastrophic) more than through its individual exponent. All
exponents are governance-adjustable parameters (see Section 8).

\textbf{Class weight rationale.} LIVENESS failures (agent crashes, API
timeouts) have the highest base recovery rate\textbackslash{}
(\textasciitilde92\% when resources are available), justifying
\emph{F}\textsubscript{LIVENESS} = 0.70. RESOURCE failures (budget
exhaustion, context overflow) are partially recoverable\textbackslash{}
(\textasciitilde48\%), justifying \emph{F}\textsubscript{RESOURCE} =
0.30. LOGIC failures (reasoning errors, hallucinations, spec mismatches)
have\textbackslash{} \textasciitilde8\% base recovery rate --- a
different agent retrying the same reasoning task rarely succeeds.
Setting \emph{F}\textsubscript{LOGIC} = 0.00 routes all LOGIC failures
directly to dispute, which is the economically correct decision: the
expected value of a recovery attempt (8\% × escrow saved) is less than
the expected cost (92\% × wasted fallback budget).

\textbf{Threshold rationale.} The Bayes-optimal three-tier sweep (Exp
13) identified \texttt{(upper,\ lower)\ =\ (0.50,\ 0.45)} as the
threshold pair that minimises overall misrouting against the
ground-truth probability \emph{p}. CAIRN ships \texttt{(0.40,\ 0.35)}
instead. The deviation is deliberate, not an error:

\begin{itemize}
\tightlist
\item
  \textbf{Asymmetric cost.} The Bayes-optimal objective treats false
  positives and false negatives symmetrically --- every misroute counts
  equally. In production economics they do not. A false positive
  (recovery attempted that fails) wastes\textbackslash{}
  \textasciitilde50\% of the remaining escrow plus the fallback agent's
  gas and reputation; a false negative (recoverable task disputed) costs
  the operator only the 3\% arbiter fee plus a 7-day delay. Section 6.6
  quantifies the asymmetry: at \emph{E} = 0.01 ETH, a FULL-tier FP costs
  \textasciitilde10× more than an FN. The Bayes-optimal threshold treats
  these as equivalent; CAIRN's lower thresholds shift the routing band
  toward the \emph{cheaper} error mode (more FNs, fewer FPs), which the
  §6.6 confusion matrix confirms (Eq4 FULL-FP drops to 7.9\% vs Bayes's
  9.2\%, FN rises to 12.2\% vs Bayes's 13.0\%).
\item
  \textbf{Coverage objective.} The narrow band \texttt{{[}0.35,\ 0.40)}
  is small by design: the multiplicative formula's primary value is the
  binary recover/dispute decision, not the FULL/REDUCED distinction.
  Setting both thresholds near the Bayes-optimal \emph{lower} boundary
  (0.45) maximises the set of tasks routed to \emph{some} recovery
  attempt rather than to immediate dispute, which is the
  operator-friendly bias.
\item
  \textbf{Headroom check.} Worked example: LIVENESS at \emph{B}=0.85,
  \emph{D}=0.88 yields \emph{r} = 0.697, comfortably above 0.40.
  RESOURCE failures span 0.25-0.45 depending on resources, sitting
  exactly in the discriminating band. LOGIC failures score 0 regardless.
\end{itemize}

All thresholds are governance-adjustable parameters (Section 8) --- an
operator population with different cost ratios can move toward (0.50,
0.45) for symmetric-cost optimisation, or further down for even more
aggressive recovery bias.

\textbf{Solidity implementation.} The multiplicative formula requires
fixed-point exponentiation. Since only three failure class weights exist
(0.70, 0.30, 0.00), the\textbackslash{} \emph{F}\^{}\emph{a} term is
implemented as a lookup (3 pre-computed values), eliminating the most
expensive \texttt{pow} call. The remaining \emph{B}\^{}\emph{b}
and\textbackslash{} \emph{D}\^{}\emph{c} terms use PRBMath's
\texttt{pow} function\textbackslash{} (\textasciitilde3,000 gas each) or
a binned lookup table for further gas savings. Total cost:
\textasciitilde2,500-6,200 gas depending on implementation ---
negligible on Base L2 (see Section 6.5).

\subsection{6.5 Gas Costs}\label{gas-costs}

All operations are designed for Base L2, where gas is inexpensive. The
\texttt{recoveryScore} and \texttt{classifyAndScore} rows below are
\textbf{measured} values from \texttt{forge\ test\ -\/-gas-report}
against the v2 reference implementation \texttt{RecoveryRouterV2.sol}
(24-test suite, 339 tests overall passing on this branch); the remaining
rows are design-target estimates pending a full v2 system benchmark. The
raw report is committed at \texttt{contracts/gas-report-v2-router.txt}.

{\def\LTcaptype{none} % do not increment counter
\begin{longtable}[]{@{}
  >{\raggedright\arraybackslash}p{(\linewidth - 8\tabcolsep) * \real{0.1594}}
  >{\raggedright\arraybackslash}p{(\linewidth - 8\tabcolsep) * \real{0.0725}}
  >{\raggedright\arraybackslash}p{(\linewidth - 8\tabcolsep) * \real{0.3768}}
  >{\raggedright\arraybackslash}p{(\linewidth - 8\tabcolsep) * \real{0.2754}}
  >{\raggedright\arraybackslash}p{(\linewidth - 8\tabcolsep) * \real{0.1159}}@{}}
\toprule\noalign{}
\begin{minipage}[b]{\linewidth}\raggedright
Operation
\end{minipage} & \begin{minipage}[b]{\linewidth}\raggedright
Gas
\end{minipage} & \begin{minipage}[b]{\linewidth}\raggedright
Cost @ 0.01 gwei Base L2
\end{minipage} & \begin{minipage}[b]{\linewidth}\raggedright
Cost @ \$2,500/ETH
\end{minipage} & \begin{minipage}[b]{\linewidth}\raggedright
Source
\end{minipage} \\
\midrule\noalign{}
\endhead
\bottomrule\noalign{}
\endlastfoot
\texttt{submitTask} & \textbackslash{} \textasciitilde180,000 (estimate)
& 1.8 × 10⁻⁶ ETH & \$0.0045 & design target \\
\texttt{commitCheckpointBatch} (1 checkpoint, v2) &
\textasciitilde80,000 (estimate) & 8.0 × 10⁻⁷ ETH & \$0.0020 & design
target \\
\texttt{commitCheckpointBatch} (10 checkpoints, v2) & \textbackslash{}
\textasciitilde100,000 (estimate) & 1.0 × 10⁻⁶ ETH & \$0.0025 & design
target \\
\texttt{heartbeat} & \textasciitilde45,000 (estimate) & 4.5 × 10⁻⁷ ETH &
\$0.0011 & design target \\
\texttt{settle} (escrow distribution) & \textbackslash{}
\textasciitilde140,000 (estimate) & 1.4 × 10⁻⁶ ETH & \$0.0035 & design
target \\
\texttt{RecoveryRouterV2.computeRecoveryScore} (full multiplicative
path) & \textbf{19,935 max / 5,748 avg / 524 min} (measured) & 2.0 ×
10⁻⁷ ETH max & \$0.00050 max & measured \\
\texttt{RecoveryRouterV2.classifyAndScore} (called from CairnCore on
failure) & \textbf{53,680 max / 39,017 avg / 24,354 min} (measured) &
5.4 × 10⁻⁷ ETH max & \$0.00134 max & measured \\
\texttt{RecoveryRouterV2} deployment cost & 1,224,782 (measured) & 1.2 ×
10⁻⁵ ETH & \$0.031 & measured \\
\end{longtable}
}

The 0.01 gwei assumption reflects typical post-Dencun Base L2 gas
prices; actual L2 execution gas has ranged from below 0.001 gwei (low
congestion) to approximately 0.1 gwei (high congestion) per BaseScan.
Base transactions also carry an L1 publication fee (\textasciitilde1-5\%
of total cost at typical congestion) that is not included in the table
above and can dominate at very low L2 gas prices. Dollar figures should
therefore be read as order-of-magnitude estimates, not contractual
guarantees.

The measured v2 \texttt{computeRecoveryScore} average of \textbf{5,748
gas} is \emph{better} than the design-target estimate\textbackslash{}
(\textasciitilde6,200 gas) used in earlier drafts. The maximum observed
(\textasciitilde20,000 gas) corresponds to the full multiplicative path
with two PRBMath UD60x18 \texttt{pow} calls; the median\textbackslash{}
(\textasciitilde1,348 gas) is dominated by short-circuit paths (LOGIC
class returning 0, or zero-budget early termination --- see
\texttt{\_computeScore} in \texttt{RecoveryRouterV2.sol}). The minimum
(524 gas) is the LOGIC short-circuit alone. The pre-computed
\emph{F}\^{}0.80 lookup (0.7518 for LIVENESS, 0.3817 for RESOURCE, 0 for
LOGIC, all at 18-decimal fixed-point precision) eliminates one PRBMath
\texttt{pow} call per score; the remaining two \texttt{pow} calls
(for\textbackslash{} \emph{B}\^{}0.35 and \emph{D}\^{}0.15) account for
most of the gas.

PRBMath UD60x18 v4.1.1 is integrated in v2
(\texttt{@prb/math/UD60x18.sol}); v1 has no fixed-point math dependency.

Merkle batching in the v2
\texttt{commitCheckpointBatch(taskId,\ count,\ merkleRoot,\ latestCID)}
function reduces checkpoint gas by approximately 95\% compared to
per-CID storage. A 50-checkpoint task is expected to cost approximately
100,000 gas for a single batch commit versus an estimated 3,350,000 gas
for 50 sequential commits at \textasciitilde67,000 gas each (linear
extrapolation; to be validated by benchmark). The v1 MVP contract uses a
simpler non-batched \texttt{commitCheckpoint(taskId,\ cid)} and does not
realize this reduction; batching is a v2 feature.

\subsection{6.6 Economic Impact of
Misrouting}\label{economic-impact-of-misrouting}

The multiplicative formula's 23.46\% misrouting rate (Section 6.4, Run 4
/ Exp 14) has a concrete economic cost derived from the Run 4 confusion
matrix (Section 10.1, Experiment 14). All percentages in this section
are \textbf{joint probabilities} P(Routed = \emph{x} ∧ Outcome =
\emph{y}) expressed over the full event population --- \emph{not}
conditional rates such as P(Failed \textbar{} Routed to Recover). The
joint-probability basis is what matters for aggregate cost; the
conditional FP/FN rates (31.10\% / 19.09\% for Eq4) appear elsewhere in
the simulation outputs and reflect a different denominator.

{\def\LTcaptype{none} % do not increment counter
\begin{longtable}[]{@{}
  >{\raggedright\arraybackslash}p{(\linewidth - 8\tabcolsep) * \real{0.2000}}
  >{\raggedright\arraybackslash}p{(\linewidth - 8\tabcolsep) * \real{0.2000}}
  >{\raggedright\arraybackslash}p{(\linewidth - 8\tabcolsep) * \real{0.2000}}
  >{\raggedright\arraybackslash}p{(\linewidth - 8\tabcolsep) * \real{0.2000}}
  >{\raggedright\arraybackslash}p{(\linewidth - 8\tabcolsep) * \real{0.2000}}@{}}
\toprule\noalign{}
\begin{minipage}[b]{\linewidth}\raggedright
Error Type
\end{minipage} & \begin{minipage}[b]{\linewidth}\raggedright
Rate (Eq4)
\end{minipage} & \begin{minipage}[b]{\linewidth}\raggedright
Rate (Eq1 Linear)
\end{minipage} & \begin{minipage}[b]{\linewidth}\raggedright
Cost per Error
\end{minipage} & \begin{minipage}[b]{\linewidth}\raggedright
Cost per 1,000 Tasks (Eq4)
\end{minipage} \\
\midrule\noalign{}
\endhead
\bottomrule\noalign{}
\endlastfoot
False positive --- FULL tier (full budget wasted) & 7.9\% & 22.3\% &
\textbackslash{} \textasciitilde50\% of remaining escrow & 79 × 0.50 ×
\emph{E}\textsubscript{rem} \\
False positive --- REDUCED tier (capped budget wasted) & 5.4\% & 3.8\% &
\textbackslash{} \textasciitilde25\% of remaining escrow & 54 × 0.25 ×
\emph{E}\textsubscript{rem} \\
False negative --- DISPUTED-but-recoverable & 12.2\% & 7.7\% & Arbiter
fee (3\% of escrow) + 7-day delay & 122 × 0.03 × \emph{E} \\
\end{longtable}
}

At an average escrow \emph{E} = 0.01 ETH with 50\%
remaining-budget\textbackslash{} \emph{E}\textsubscript{rem} = 0.005 ETH
at failure (so per 1,000 tasks, each percentage-point of joint
probability corresponds to 10 events):

\begin{itemize}
\tightlist
\item
  FULL-tier FP: 79 × 0.50 × 0.005 ETH = \textbf{0.198 ETH / 1,000 tasks}
\item
  REDUCED-tier FP: 54 × 0.25 × 0.005 ETH = \textbf{0.068 ETH / 1,000
  tasks}
\item
  FN (arbiter fee only, delay cost not monetized): 122 × 0.03 × 0.01 ETH
  = \textbf{0.037 ETH / 1,000 tasks}
\item
  \textbf{Total Eq4 misrouting cost:\textbackslash{}
  \textasciitilde0.303 ETH / 1,000 tasks (\textasciitilde\$758 at
  \$2,500/ETH, i.e.,\textbackslash{} \textasciitilde\$0.76 per task).}
\end{itemize}

\textbf{Versus the current Eq1 formula (v1 testnet deployment).}
Applying the same cost model to the Run 1 Eq1-current confusion matrix
(RESULTS.md §9: FULL+Failed 13.1\%, REDUCED+Failed 34.0\%,
DISPUTED+Succeeded 0.5\%):

\begin{itemize}
\tightlist
\item
  FULL-tier FP: 131 × 0.50 × 0.005 ETH = 0.3275 ETH / 1,000 tasks
\item
  REDUCED-tier FP: 340 × 0.25 × 0.005 ETH = 0.4250 ETH / 1,000 tasks
\item
  FN: 5 × 0.03 × 0.01 ETH = 0.0015 ETH / 1,000 tasks
\item
  \textbf{Total Eq1 misrouting cost: \textasciitilde0.754 ETH / 1,000
  tasks\textbackslash{} (\textasciitilde\$1,885).}
\end{itemize}

The multiplicative v2 formula therefore saves
\textbf{\textasciitilde0.451 ETH / 1,000 tasks\textbackslash{}
(\textasciitilde\$1,128)}, a \textbf{60\% reduction} in misrouting cost.
The residual 0.93pp gap above the Bayes-optimal floor (22.53\%) bounds
the maximum further savings from calibration at roughly 0.012 ETH /
1,000 tasks (\textasciitilde\$30) --- essentially exhausted.

Relative to deployed escrow capital (10 ETH of escrow across 1,000 tasks
at \emph{E} = 0.01): Eq4 misrouting costs\textbackslash{}
\textbf{\textasciitilde3.0\% of escrow value}; Eq1-current costs
\textasciitilde7.5\%. Eq4 reduces the friction on deployed escrow
by\textbackslash{} \textasciitilde2.5× --- an acceptable overhead for
permissionless, trustless automated recovery.

\begin{center}\rule{0.5\linewidth}{0.5pt}\end{center}

\section{7. Security Model}\label{security-model}

\subsection{7.1 Design Principles}\label{design-principles}

{\def\LTcaptype{none} % do not increment counter
\begin{longtable}[]{@{}
  >{\raggedright\arraybackslash}p{(\linewidth - 2\tabcolsep) * \real{0.4231}}
  >{\raggedright\arraybackslash}p{(\linewidth - 2\tabcolsep) * \real{0.5769}}@{}}
\toprule\noalign{}
\begin{minipage}[b]{\linewidth}\raggedright
Principle
\end{minipage} & \begin{minipage}[b]{\linewidth}\raggedright
Implementation
\end{minipage} \\
\midrule\noalign{}
\endhead
\bottomrule\noalign{}
\endlastfoot
No trusted keepers & All enforce functions public; anyone can trigger
failure detection \\
Escrow as commitment & Record-writing mandatory for payment release \\
Stake-based accountability & Fallback agents and arbiters stake capital
proportional to exposure \\
Deterministic scoring & Recovery score is a pure function of on-chain
state --- no oracle, no AI \\
Permissionless enforcement & False enforcement calls revert; caller pays
gas; no economic benefit to attacker \\
\end{longtable}
}

\subsection{7.2 Threat Model}\label{threat-model}

CAIRN assumes: - The underlying blockchain (Base/Ethereum) is secure -
IPFS content remains available for the task duration plus dispute period
(ensured via pinning services) - ERC-8004 reputation scores are accurate
(CAIRN inherits ERC-8004's security model) - Operators submit accurate
task specifications (agents can query specs before accepting) - Block
time is consistent (\textasciitilde2s on Base)

\textbf{L2 sequencer trust dependency.} Base is operated by a single
centralised sequencer (Coinbase, as of April 2026). For a protocol that
markets ``trustless'' recovery, this dependency requires explicit
treatment:

\begin{itemize}
\tightlist
\item
  \textbf{Liveness:} if the Coinbase sequencer is offline or refuses to
  include CAIRN transactions, the entire protocol pauses ---
  \texttt{checkLiveness}, \texttt{commitCheckpointBatch},
  \texttt{settle}, and arbiter rulings cannot execute. This is a
  censorship/availability risk shared with every Base-deployed protocol.
  The mitigation Base provides today is a ``force inclusion'' escape
  hatch through L1, but it adds latency (hours-to-days, depending on
  Base's exact bridge cadence) that exceeds CAIRN's heartbeat intervals;
  force-included transactions would arrive too late to prevent stale
  failures.
\item
  \textbf{Ordering:} the sequencer chooses the in-block order of
  transactions. For CAIRN this matters in two cases: (i) a worker agent
  submitting a just-in-time \texttt{heartbeat} racing against an
  enforcer's \texttt{checkLiveness}, and (ii) a fallback agent
  committing checkpoints racing against the deadline. In both cases the
  sequencer can pick a winner. The atomicity of CAIRN's failure path
  (detection → classification → routing in one transaction) limits the
  MEV surface to ordering only --- there is no mid-transaction state
  insertion --- but ordering alone is sufficient to extract value in
  adversarial scenarios.
\item
  \textbf{Honesty assumption:} CAIRN implicitly assumes the sequencer is
  not actively collaborating with one party to a dispute. This
  assumption is weaker than the underlying L1 trust model and is the
  principal residual trust dependency in v2. Base's roadmap toward
  decentralised sequencing (planned via the Optimism Superchain stack)
  will reduce this dependency over time; until then, CAIRN inherits the
  same trust assumption as every other Base-deployed application.
\item
  \textbf{Migration mitigation:} the protocol is portable to any EVM L2
  with comparable gas economics. A future deployment to a chain with
  decentralised sequencing (e.g., a future iteration of the OP Stack
  with shared sequencing, or an Arbitrum Nitro chain with permissionless
  validators) would remove the single-sequencer dependency without
  protocol changes.
\end{itemize}

\subsection{7.3 Attack Vectors and
Mitigations}\label{attack-vectors-and-mitigations}

{\def\LTcaptype{none} % do not increment counter
\begin{longtable}[]{@{}
  >{\raggedright\arraybackslash}p{(\linewidth - 4\tabcolsep) * \real{0.2667}}
  >{\raggedright\arraybackslash}p{(\linewidth - 4\tabcolsep) * \real{0.3333}}
  >{\raggedright\arraybackslash}p{(\linewidth - 4\tabcolsep) * \real{0.4000}}@{}}
\toprule\noalign{}
\begin{minipage}[b]{\linewidth}\raggedright
Attack
\end{minipage} & \begin{minipage}[b]{\linewidth}\raggedright
Severity
\end{minipage} & \begin{minipage}[b]{\linewidth}\raggedright
Mitigation
\end{minipage} \\
\midrule\noalign{}
\endhead
\bottomrule\noalign{}
\endlastfoot
\textbf{Checkpoint gaming} --- Agent commits fake checkpoints to inflate
payment & High & v2: schema hash validation on commit (the
\texttt{specHash} stored at task initialization is matched against each
checkpoint payload); off-chain content verification; reputation decay
for invalid submissions. (v1 stores \texttt{specHash} but does not
enforce per-checkpoint validation --- the v2 upgrade activates this
check.) \\
\textbf{Liveness griefing} --- Attacker calls \texttt{checkLiveness} to
force premature failure & Low & Only succeeds if heartbeat interval
actually elapsed; false calls revert; attacker pays gas \\
\textbf{Fallback Sybil} --- Attacker registers many agents to capture
recovery assignments & Medium & Reputation gate (min 50) + stake
requirement (10\%); 100\% slash on zero-checkpoint failure \\
\textbf{Arbiter collusion} --- Arbiter rules in favor of colluding agent
& Medium & Stake (20\%) \textgreater{} fee (3\%); incorrect rulings
slashed; commit-reveal prevents front-running \\
\textbf{Recovery score manipulation} --- Agent times failure for desired
routing & Low & All score inputs on-chain; agent cannot control failure
classification retroactively \\
\textbf{Intelligence poisoning} --- False failure records to mislead
future agents & Medium & Records auto-written by protocol, not by
agents; content matches on-chain state \\
\textbf{Sequencer reordering (MEV)} --- Block builder reorders
\texttt{checkLiveness} and a just-in-time \texttt{heartbeat} to extract
value & Low & Atomic detection → classification → routing within a
single tx; MEV surface limited to ordering, not mid-tx state insertion;
heartbeat interval must actually have elapsed for enforcement to
succeed \\
\end{longtable}
}

\subsection{7.4 Protocol Invariants}\label{protocol-invariants}

These properties hold at all times:

\begin{enumerate}
\def\labelenumi{\arabic{enumi}.}
\tightlist
\item
  \textbf{Escrow safety.} Escrow is not released until
  \texttt{state\ ==\ RESOLVED}.
\item
  \textbf{Deterministic scoring.} Same on-chain inputs produce the same
  recovery score. No external dependencies.
\item
  \textbf{Checkpoint immutability.} Committed checkpoint CIDs cannot be
  modified or deleted.
\item
  \textbf{State irreversibility.} No backward state transitions. The DAG
  is
  \texttt{IDLE\ →\ RUNNING\ →\ \{FAILED,\ RESOLVED\};\ FAILED\ →\ \{RECOVERING,\ DISPUTED\};\ RECOVERING\ →\ \{RESOLVED,\ DISPUTED\};\ DISPUTED\ →\ RESOLVED;\ RESOLVED\ →\ terminal}.
  The \texttt{RECOVERING\ →\ DISPUTED} edge exists in the v2
  specification as a direct transition when the fallback agent itself
  fails (FallbackFailed event in §2.2.1, triggered by missed heartbeat
  or deadline exceeded during the fallback phase). In v1, this
  transition is reached \emph{indirectly} via the fallback's failure
  re-entering the FAILED state and then being re-routed; v2 surfaces it
  as a direct edge to make the state graph match the formal model in
  this paper. The reachable state set is identical in both versions; v2
  changes only the path, not the destinations.
\item
  \textbf{Fee ordering.} Protocol fee is deducted before agent payment
  calculation.
\item
  \textbf{Liveness enforcement.} \texttt{checkLiveness} only succeeds if
  the heartbeat interval has actually elapsed.
\end{enumerate}

\subsection{7.5 Incentive Analysis}\label{incentive-analysis}

We analyze three strategic games within CAIRN: the checkpoint commitment
game, the fallback acceptance game, and the arbiter ruling game.

\subsubsection{Game 1: Checkpoint
Commitment}\label{game-1-checkpoint-commitment}

\textbf{Proposition 1 (Checkpoint Dominance).} \emph{For a task with
escrow E, n subtasks, empirical failure probability\textbackslash{}
p\textsubscript{f}, and per-checkpoint gas cost g, honest checkpointing
is the dominant strategy when\textbackslash{} p\textsubscript{f} × E ×
(1 − f) / n \textgreater{} g, where f is the protocol fee rate.}

\emph{Proof.} Agent \emph{A} executes task \emph{T} with \emph{n}
subtasks. At each subtask \emph{i}, \emph{A} chooses
action\textbackslash{} \emph{a}\textsubscript{i} ∈ \{commit, skip\}. Let
\emph{c} = number of committed checkpoints. Let \emph{g} = gas cost per
checkpoint on Base L2 ≈ 8 × 10⁻⁷ ETH (80,000 gas × 0.01 gwei).

The expected marginal value of one additional checkpoint:

\begin{verbatim}
ΔV = (1 − p_f) × (−g) + p_f × [E(1 − f)/n − g]
   = p_f × E(1 − f)/n − g
\end{verbatim}

In case of success (probability 1 −\textbackslash{}
\emph{p}\textsubscript{f}), an extra checkpoint costs \emph{g} but
provides no additional payment. In case of failure
(probability\textbackslash{} \emph{p}\textsubscript{f}), an extra
checkpoint increases the agent's share by \emph{E}(1 −
\emph{f})/\emph{n} at cost \emph{g}.

Checkpointing is dominant when Δ\emph{V} \textgreater{} 0. Evaluating
with protocol parameters (\emph{f} = 0.005, \emph{g} = 8 × 10⁻⁷ ETH) and
the empirical failure rate\textbackslash{} \emph{p}\textsubscript{f} =
0.5 from {[}3{]}:

{\def\LTcaptype{none} % do not increment counter
\begin{longtable}[]{@{}llll@{}}
\toprule\noalign{}
Subtasks (\emph{n}) & Min Escrow (\emph{E}) & Δ\emph{V} & Dominant? \\
\midrule\noalign{}
\endhead
\bottomrule\noalign{}
\endlastfoot
5 & 0.001 ETH & 9.87 × 10⁻⁵ & Yes (Δ\emph{V} ≫ \emph{g}) \\
10 & 0.001 ETH & 4.90 × 10⁻⁵ & Yes \\
50 & 0.001 ETH & 9.15 × 10⁻⁶ & Yes \\
100 & 0.001 ETH & 4.17 × 10⁻⁶ & Yes \\
500 & 0.001 ETH & 1.95 × 10⁻⁷ & Marginal (Δ\emph{V} ≈ 0.24\emph{g}) \\
\end{longtable}
}

The condition holds for all realistic task configurations (n ≤ 100 at
minimum escrow, or any escrow ≥ 0.01 ETH at any \emph{n}). The boundary
escrow for dominance is\textbackslash{} \emph{E}\textsubscript{min} =
(\emph{g} × \emph{n}) /\textbackslash{} (\emph{p}\textsubscript{f} × (1
− \emph{f})) --- for \emph{n} = 100:\textbackslash{}
\emph{E}\textsubscript{min} ≈ 1.6 × 10⁻⁴ ETH, well below the protocol
minimum of 10⁻³ ETH.

\textbf{Faking checkpoints} is strictly dominated under v2: invalid CIDs
are rejected by schema hash validation (the \texttt{specHash} committed
at task initialization and matched against each checkpoint payload).
Under v1 (where per-checkpoint validation is not yet enforced), the
equilibrium is preserved by reputation decay alone --- repeated invalid
submissions push agents below the fallback pool admission threshold
(Section 4.4). The v2 upgrade hardens this from an economic disincentive
into an on-chain invariant. ∎

\subsubsection{Game 2: Fallback
Acceptance}\label{game-2-fallback-acceptance}

\textbf{Proposition 2 (Rational Fallback Acceptance).} \emph{A fallback
agent with success rate\textbackslash{} s\textsubscript{τ} for task type
τ should accept a recovery assignment when the expected payoff exceeds
zero:}

\begin{verbatim}
s_τ × E_r × (c_f / c_total) × (1 − f) > (1 − s_τ) × p_0 × stake
\end{verbatim}

\emph{where\textbackslash{} E\textsubscript{r} is remaining
escrow,\textbackslash{} c\textsubscript{f}/c\textsubscript{total} is the
fallback's expected checkpoint share, and\textbackslash{}
p\textsubscript{0} is the probability of zero-checkpoint failure
(triggering 100\% slash).}

\emph{Proof.} A fallback agent accepting an assignment faces a binary
outcome: - With probability\textbackslash{} \emph{s}\textsubscript{τ}
(success), the fallback earns its share of remaining escrow after the
protocol fee: gain =\textbackslash{} \emph{E}\textsubscript{r}
×\textbackslash{} (\emph{c}\textsubscript{f} /\textbackslash{}
\emph{c}\textsubscript{total}) × (1 − \emph{f}). - With probability (1
−\textbackslash{} \emph{s}\textsubscript{τ}) (failure), the fallback
either (i) completes at least one checkpoint and earns partial payment
without stake loss, or (ii) fails without any checkpoint and is slashed
100\% of stake.

Let\textbackslash{} \emph{p}\textsubscript{0} denote the conditional
probability of zero-checkpoint failure given failure --- the sole slash
trigger. Expected payoff:

\begin{verbatim}
E[payoff] = s_τ · E_r · (c_f / c_total) · (1 − f)  −  (1 − s_τ) · p_0 · stake
\end{verbatim}

Acceptance is rational when E{[}payoff{]} \textgreater{} 0, yielding the
stated inequality. Three remarks: (a) at\textbackslash{}
\emph{s}\textsubscript{τ} = 1 the condition reduces to a trivially
positive gain, recovering the intuition that a perfectly competent
fallback always accepts; (b) the fallback's own gas cost \emph{g} and
opportunity cost are absorbed into\textbackslash{}
\emph{E}\textsubscript{r} for brevity (subtract them from the left side
to obtain the stricter form); (c) the admission gates (reputation ≥ 50,
stake ≥ 10\%) enforce a lower bound on\textbackslash{}
\emph{s}\textsubscript{τ} by construction, so low-competence agents
never reach the decision. ∎

This self-selection mechanism is beneficial: agents with low success
rates for a given task type will rationally decline recovery
assignments, ensuring only capable agents accept.

\subsubsection{Game 3: Arbiter Ruling}\label{game-3-arbiter-ruling}

\textbf{Proposition 3 (Arbiter Honesty).} \emph{For an arbiter with
stake s = 0.2V (where V is the dispute value), fee φ = 0.03V, and
detection probability\textbackslash{} p\textsubscript{d} for incorrect
rulings, honest ruling is the rational strategy when:}

\begin{verbatim}
bribe < φ + p_d × 0.5s = 0.03V + p_d × 0.1V
\end{verbatim}

\emph{Proof.} The arbiter chooses between honest and dishonest ruling: -
\textbf{Honest payoff:} +\emph{φ} = +0.03\emph{V} - \textbf{Dishonest
payoff:} +bribe −\textbackslash{} \emph{p}\textsubscript{d} ×
0.5\emph{s} = bribe −\textbackslash{} \emph{p}\textsubscript{d} ×
0.1\emph{V}

For dishonesty to be rational, the bribe must exceed the honest fee plus
the expected penalty: bribe \textgreater{} 0.03\emph{V}
+\textbackslash{} \emph{p}\textsubscript{d} × 0.1\emph{V}.

The detection probability\textbackslash{} \emph{p}\textsubscript{d}
varies sharply by dispute type, and this is the central weakness of the
analysis. For \textbf{LIVENESS} disputes (was a heartbeat actually
missed?) and \textbf{RESOURCE} disputes (was the budget actually
exhausted?), all evidence is on-chain --- heartbeat timestamps, cost
accruals, deadline blocks are publicly verifiable --- so\textbackslash{}
\emph{p}\textsubscript{d} ≥ 0.8 is conservative. With\textbackslash{}
\emph{p}\textsubscript{d} = 0.8 the minimum bribe is 0.03\emph{V} +
0.08\emph{V} = 0.11\emph{V}.

\textbf{Multi-dispute analysis:} For LIVENESS/RESOURCE classes, an
arbiter who colludes in \emph{k} disputes has cumulative detection
probability 1 − (1 −\textbackslash\textbackslash{}
\emph{p}\textsubscript{d})\textsuperscript{k}. At\textbackslash{}
\emph{p}\textsubscript{d} = 0.8: a single dispute has 80\% detection
risk; two disputes have 96\%; three disputes have 99.2\%. Sustained
collusion is therefore economically irrational for these classes ---
expected losses compound exponentially while gains are linear.

\textbf{The arbiter-recursion problem.} Proposition 3's economic
guarantee breaks down for \textbf{LOGIC} disputes (``did the agent
hallucinate?'', ``did the output match the specification?''). Arbiters
are themselves agents reasoning about reasoning errors; the dispute
being judged and the judge are drawn from the same population. Whether a
ruling is ``incorrect'' is not always detectable from on-chain evidence
--- for genuine LOGIC disputes the detection signal is itself an
inference about reasoning quality, not a comparison against a timestamp.
We estimate\textbackslash{} \emph{p}\textsubscript{d} for the hardest
LOGIC cases at 0.3-0.5 rather than 0.8, which raises the minimum bribe
to roughly 0.045\emph{V}-0.08\emph{V} --- still meaningful, but a much
narrower margin than the on-chain-evidenced classes.

This is the same structural problem encountered by Kleros and other
on-chain dispute-resolution protocols {[}19{]}: a sufficiently capable
adversary can craft a dispute where the ``correct'' answer is genuinely
contestable, and detection requires either (a) an appeals layer that
defers to ever-larger juries with diminishing returns, or (b) an
external ground-truth oracle that re-introduces a trust assumption.
CAIRN inherits this limitation. The mitigations the protocol does
provide --- high reputation gates for arbiter admission, commit-reveal
to prevent pre-arrangement, multi-dispute compounding for repeat
offenders --- bound the worst case but do not eliminate it. Section 10.3
lists this as a known limitation; the v2 specification ships
\textbf{single-tier arbitration}, and a future appeals layer
(multi-juror with bonded escalation, in the Kleros style) is reserved
for protocol v3.

The commit-reveal scheme provides an additional layer: the arbiter
commits a hash of their ruling before knowing which agents are involved
(preventing pre-arrangement), and must reveal within 24 hours or forfeit
their eligibility. ∎

\begin{center}\rule{0.5\linewidth}{0.5pt}\end{center}

\section{8. Governance}\label{governance}

\subsection{8.1 Configurable Parameters}\label{configurable-parameters}

{\def\LTcaptype{none} % do not increment counter
\begin{longtable}[]{@{}
  >{\raggedright\arraybackslash}p{(\linewidth - 6\tabcolsep) * \real{0.2750}}
  >{\raggedright\arraybackslash}p{(\linewidth - 6\tabcolsep) * \real{0.2250}}
  >{\raggedright\arraybackslash}p{(\linewidth - 6\tabcolsep) * \real{0.1750}}
  >{\raggedright\arraybackslash}p{(\linewidth - 6\tabcolsep) * \real{0.3250}}@{}}
\toprule\noalign{}
\begin{minipage}[b]{\linewidth}\raggedright
Parameter
\end{minipage} & \begin{minipage}[b]{\linewidth}\raggedright
Default
\end{minipage} & \begin{minipage}[b]{\linewidth}\raggedright
Range
\end{minipage} & \begin{minipage}[b]{\linewidth}\raggedright
Description
\end{minipage} \\
\midrule\noalign{}
\endhead
\bottomrule\noalign{}
\endlastfoot
Protocol fee & 50 bps & 0-500 bps & Fee on escrow settlement \\
Fallback min reputation & 50/100 & 0-100 & Min ERC-8004 reputation for
fallback pool \\
Fallback min stake & 10\% & 1-50\% & Stake as \% of max eligible
escrow \\
Arbiter min stake & 20\% & 5-50\% & Stake as \% of max ruleable
dispute \\
Arbiter fee & 3\% & 1-10\% & Fee as \% of dispute value \\
Arbiter timeout & 7 days & 1-30 days & Time for arbiter to rule before
auto-refund \\
Recovery threshold (upper) & 0.40 & 0.1-0.9 & Score threshold for
full-scope RECOVERING \\
Recovery threshold (lower) & 0.35 & 0.1-0.9 & Score threshold for
reduced-scope RECOVERING vs DISPUTED \\
Recovery exponent \emph{a} & 0.80 & 0.1-2.0 & Exponent for failure class
weight in Equation 1 \\
Recovery exponent \emph{b} & 0.35 & 0.1-2.0 & Exponent for budget
remaining in Equation 1 \\
Recovery exponent \emph{c} & 0.15 & 0.1-2.0 & Exponent for deadline
remaining in Equation 1 \\
\textbackslash{} \emph{F}\textsubscript{LIVENESS} & 0.70 & 0-1.0 &
Failure class weight for LIVENESS failures \\
\textbackslash{} \emph{F}\textsubscript{RESOURCE} & 0.30 & 0-1.0 &
Failure class weight for RESOURCE failures \\
\textbackslash{} \emph{F}\textsubscript{LOGIC} & 0.00 & 0-1.0 & Failure
class weight for LOGIC failures \\
\end{longtable}
}

\subsection{8.2 Governance Model}\label{governance-model}

\textbf{Current: Multi-sig with timelock.} A 3-of-5 multi-sig controls
parameter changes with a 48-hour timelock. All proposals are public.
This provides rapid iteration capability while maintaining review
periods.

\textbf{Future: Token governance (if applicable).} Parameters changeable
via on-chain voting. Emergency multi-sig retained for security issues.

\subsection{8.3 Upgrade Path}\label{upgrade-path}

CAIRN contracts implement the \textbf{UUPS (Universal Upgradeable Proxy
Standard)} pattern. The proxy is immutable; the implementation contract
can be replaced via governance-authorized upgrade. This allows the
protocol to evolve (new features, bug fixes) while preserving user funds
and task state across versions.

All upgrades require governance approval via the timelock. In-flight
tasks are not affected --- they complete under the implementation
version active when they were created.

\begin{center}\rule{0.5\linewidth}{0.5pt}\end{center}

\section{9. Related Work}\label{related-work}

\subsection{9.1 Agent Frameworks with Failure
Handling}\label{agent-frameworks-with-failure-handling}

{\def\LTcaptype{none} % do not increment counter
\begin{longtable}[]{@{}
  >{\raggedright\arraybackslash}p{(\linewidth - 8\tabcolsep) * \real{0.3333}}
  >{\centering\arraybackslash}p{(\linewidth - 8\tabcolsep) * \real{0.1667}}
  >{\centering\arraybackslash}p{(\linewidth - 8\tabcolsep) * \real{0.1667}}
  >{\centering\arraybackslash}p{(\linewidth - 8\tabcolsep) * \real{0.1667}}
  >{\centering\arraybackslash}p{(\linewidth - 8\tabcolsep) * \real{0.1667}}@{}}
\toprule\noalign{}
\begin{minipage}[b]{\linewidth}\raggedright
Solution
\end{minipage} & \begin{minipage}[b]{\linewidth}\centering
Checkpointing
\end{minipage} & \begin{minipage}[b]{\linewidth}\centering
Cross-Agent Handoff
\end{minipage} & \begin{minipage}[b]{\linewidth}\centering
Escrow Settlement
\end{minipage} & \begin{minipage}[b]{\linewidth}\centering
Failure Intelligence
\end{minipage} \\
\midrule\noalign{}
\endhead
\bottomrule\noalign{}
\endlastfoot
LangGraph & Yes (proprietary) & No & No & No \\
Temporal.io & Yes & No (same worker) & No & No \\
Kubernetes & No (container-level) & No & No & No \\
LangSmith & No (observability only) & No & No & No \\
Manual restart & No & No & No & No \\
\textbf{CAIRN} & \textbf{Yes (portable)} & \textbf{Yes} & \textbf{Yes} &
\textbf{Yes} \\
\end{longtable}
}

\textbf{LangGraph} provides state persistence via
\texttt{SqliteSaver}/\texttt{PostgresSaver}, but checkpoints are
framework-locked --- a LangGraph checkpoint cannot be read by a CrewAI
agent or an Olas mech. No automatic fallback assignment, no escrow
settlement, no cross-ecosystem intelligence. CAIRN wraps LangGraph:
agents write CAIRN-format checkpoints (portable via IPFS), and any
CAIRN-compatible agent can resume on failure.

\textbf{Temporal.io} provides production-grade workflow orchestration
with deterministic replay. However, recovery is within the same worker
pool --- no cross-organization fallback. No escrow integration. No
on-chain settlement. Temporal is excellent for internal automation;
CAIRN is for autonomous agents with on-chain economics.

\textbf{Kubernetes} operates at the container level (is the pod
running?) while CAIRN operates at the task level (did step 3 complete?
can step 4 be recovered?). These are orthogonal. Use both.

\subsection{9.2 On-Chain Agent
Coordination}\label{on-chain-agent-coordination}

\textbf{ERC-8183 / Agent Commerce Protocol (ACP).} Defines the job
lifecycle (Open → Funded → Submitted → Terminal) and escrow mechanism.
Live on Arbitrum via Virtuals Protocol {[}15{]}. CAIRN extends ERC-8183
as a Hook --- when an ERC-8183 job fails, CAIRN provides the failure
classification, fallback routing, and proportional settlement that
ERC-8183 does not define.

\textbf{ERC-8004 (Trustless Agents).} Provides agent identity and
reputation registries. Live on Ethereum mainnet since January 29, 2026,
with\textbackslash{} \textasciitilde49,000 registered agents across 30+
EVM chains as of February 2026 {[}14{]}. CAIRN reads reputation scores
for fallback pool admission and writes outcome signals (success/failure
attestations) after task completion.

\textbf{ERC-7710 (Delegation Framework).} Enables scoped permission
transfer. CAIRN uses ERC-7710 for pre-authorized fallback delegation:
the operator grants CAIRN permission to sub-delegate task authority to a
fallback agent at recovery time, without requiring a new operator
signature.

\textbf{ERC-8211 (Smart Batching).} Enables AI agents to execute complex
DeFi operations in a single batched transaction {[}16{]}. Composable
with CAIRN: checkpoints can wrap ERC-8211 batched executions.

\subsection{9.3 Off-Chain Coordination
Protocols}\label{off-chain-coordination-protocols}

\textbf{Google A2A (Agent-to-Agent Protocol).} Donated to the Linux
Foundation (June 2025) with 100+ industry partners {[}12{]}. Handles
agent discovery via Agent Cards, task delegation, and collaboration via
context sharing. A2A is the off-chain communication layer; CAIRN is the
on-chain settlement and recovery layer. A2A tells agents what to do;
CAIRN guarantees what happens when it fails.

\textbf{Anthropic MCP (Model Context Protocol).} Connects agents to
tools and data sources {[}13{]}. Adopted by OpenAI, Microsoft, Google,
Amazon. MCP provides the tool access; CAIRN provides the fault
tolerance. When an MCP tool call fails mid-task, CAIRN's checkpoint
system preserves completed work and the fallback mechanism ensures
completion.

\subsection{9.4 Academic Foundations}\label{academic-foundations}

CAIRN builds on established theory:

\begin{itemize}
\item
  \textbf{Distributed checkpointing.} The Chandy-Lamport algorithm
  {[}8{]} proves that consistent global state can be reconstructed from
  local checkpoints. CAIRN adapts this for AI agent semantics with
  IPFS-stored, schema-validated checkpoints.
\item
  \textbf{Mechanism design.} Staking and slashing mechanisms are proven
  at over \$100 billion in staked value in Ethereum Proof-of-Stake
  {[}9{]}. CAIRN applies the same principle: stake capital to
  participate, lose it for misbehavior.
\item
  \textbf{Multi-agent failure analysis.} The MAST taxonomy {[}1{]}
  provides the most comprehensive classification of multi-agent failure
  modes to date. CAIRN's contribution is to add the recoverability
  dimension --- classifying failures by what action to take, not by what
  went wrong.
\end{itemize}

\begin{center}\rule{0.5\linewidth}{0.5pt}\end{center}

\section{10. Future Work}\label{future-work}

\subsection{10.1 Open Research Questions}\label{open-research-questions}

\textbf{Checkpoint semantic portability.} CAIRN's checkpoint format is
portable across frameworks (schema-validated IPFS CIDs), and Section
4.1.1 establishes that fully portable and portable-with-context
checkpoints cover the majority of practical task types (data fetches,
API calls, multi-step computations, stateful queries). However, for
complex reasoning chains with implicit context (chain-of-thought,
multi-turn dialogue), semantic portability remains unproven. Empirical
study of checkpoint portability across LangGraph, CrewAI, and Olas agent
architectures is planned.

\textbf{Recovery score calibration --- empirical validation.} The
recovery score formula (Section 6.4) has been validated via a Monte
Carlo simulation suite of 100,000 synthetic task-failure events per run,
seed=42, fully reproducible (Reference {[}18{]}). The simulation
compared four formula structures across \textbf{16 experiments organized
into 4 runs}. Run 1 evaluated 362 configurations (55 weight vectors +
245 class-weight vectors + 62 threshold pairs, sequential grid
refinement); Run 4 evaluated 2,646 multiplicative-formula
configurations; Runs 2-3 used staged grid refinement within the extended
parameter space (counts emit to \texttt{stdout} at runtime).

\emph{Run summary.} Each run tested a distinct formula hypothesis:

{\def\LTcaptype{none} % do not increment counter
\begin{longtable}[]{@{}
  >{\raggedright\arraybackslash}p{(\linewidth - 8\tabcolsep) * \real{0.0962}}
  >{\raggedright\arraybackslash}p{(\linewidth - 8\tabcolsep) * \real{0.1731}}
  >{\raggedright\arraybackslash}p{(\linewidth - 8\tabcolsep) * \real{0.2115}}
  >{\raggedright\arraybackslash}p{(\linewidth - 8\tabcolsep) * \real{0.2115}}
  >{\raggedright\arraybackslash}p{(\linewidth - 8\tabcolsep) * \real{0.3077}}@{}}
\toprule\noalign{}
\begin{minipage}[b]{\linewidth}\raggedright
Run
\end{minipage} & \begin{minipage}[b]{\linewidth}\raggedright
Formula
\end{minipage} & \begin{minipage}[b]{\linewidth}\raggedright
Structure
\end{minipage} & \begin{minipage}[b]{\linewidth}\raggedright
Variables
\end{minipage} & \begin{minipage}[b]{\linewidth}\raggedright
Misrouting Rate
\end{minipage} \\
\midrule\noalign{}
\endhead
\bottomrule\noalign{}
\endlastfoot
1 & \emph{r} = \emph{w}\textsubscript{f}\emph{F} +\textbackslash{}
\emph{w}\textsubscript{b}\emph{B} +\textbackslash{}
\emph{w}\textsubscript{d}\emph{D} & Linear & 3 (\emph{F}, \emph{B},
\emph{D}) & 47.56\% (current) → \textbf{33.81\%} (optimized) \\
2 & Run 1 + piecewise cliffs +\textbackslash{}
\emph{w}\textsubscript{int}·\emph{B}·\emph{D} interaction & Linear +
non-linear & 3 + 4 cliff params + 1 interaction & \textbf{33.17\%} \\
3 & Run 1 +\textbackslash{} \emph{w}\textsubscript{c}·\emph{C}
(complexity) +\textbackslash{} \emph{w}\textsubscript{s}·\emph{S}
(fallback skill) & Linear & 5 (adds \emph{C}, \emph{S}) &
\textbf{32.78\%} \\
\textbf{4} & \textbf{\emph{r} =\textbackslash{} \emph{F}\^{}\emph{a} ×
\emph{B}\^{}\emph{b} ×\textbackslash{} \emph{D}\^{}\emph{c}} &
\textbf{Multiplicative} & \textbf{3 (\emph{F}, \emph{B}, \emph{D})} &
\textbf{23.46\%} \\
--- & Bayes-optimal three-tier (thresholds 0.50/0.45) & Perfect oracle
(ground-truth \emph{p}) & --- & 22.53\% \\
\end{longtable}
}

Runs 1-3 exhaustively proved that any additive formula --- regardless of
non-linear terms or additional variables --- converges to
a\textbackslash{} \textasciitilde33\% misrouting structural ceiling. The
\textbf{``93/4/3 rule''} emerged: 93\% of achievable improvement comes
from weight tuning within the linear formula, 4\% from non-linear terms,
and 3\% from adding variables. The ceiling exists because additive
formulas cannot express the ``any-factor-kills-it'' dynamic: in reality,
zero budget means zero recovery chance regardless of failure type, but a
sum always produces a positive value from the remaining terms.

Run 4 changed the formula structure to multiplicative. The result ---
23.46\% misrouting --- captures 96\% of the theoretically achievable
improvement and lies within 0.93 percentage points of the Bayes-optimal
three-tier baseline (22.53\%). The binary Bayes-optimal floor is
22.52\%, confirming that the 22.5\% irreducible noise is intrinsic to
the stochastic ground truth rather than to the three-tier structure.

\emph{Experiment catalog.} Each experiment answered a distinct
calibration question:

{\def\LTcaptype{none} % do not increment counter
\begin{longtable}[]{@{}
  >{\raggedright\arraybackslash}p{(\linewidth - 10\tabcolsep) * \real{0.1136}}
  >{\raggedright\arraybackslash}p{(\linewidth - 10\tabcolsep) * \real{0.1136}}
  >{\raggedright\arraybackslash}p{(\linewidth - 10\tabcolsep) * \real{0.1364}}
  >{\raggedright\arraybackslash}p{(\linewidth - 10\tabcolsep) * \real{0.1818}}
  >{\raggedright\arraybackslash}p{(\linewidth - 10\tabcolsep) * \real{0.2955}}
  >{\raggedright\arraybackslash}p{(\linewidth - 10\tabcolsep) * \real{0.1591}}@{}}
\toprule\noalign{}
\begin{minipage}[b]{\linewidth}\raggedright
Exp
\end{minipage} & \begin{minipage}[b]{\linewidth}\raggedright
Run
\end{minipage} & \begin{minipage}[b]{\linewidth}\raggedright
Name
\end{minipage} & \begin{minipage}[b]{\linewidth}\raggedright
Method
\end{minipage} & \begin{minipage}[b]{\linewidth}\raggedright
Key Finding
\end{minipage} & \begin{minipage}[b]{\linewidth}\raggedright
Delta
\end{minipage} \\
\midrule\noalign{}
\endhead
\bottomrule\noalign{}
\endlastfoot
1 & 1 & Weight optimization & Grid search over 55
\emph{(w\textsubscript{f},\textbackslash{}
w\textsubscript{b},\textbackslash{} w\textsubscript{d})} simplex points
& Current weights rank \#31/55; deadline should dominate\textbackslash{}
(\emph{w\textsubscript{d}} ≥ 0.40) & −5.35pp \\
2 & 1 & Class weight optimization & Grid search over 245\textbackslash{}
\emph{(F\textsubscript{L},\textbackslash{}
F\textsubscript{R},\textbackslash{} F\textsubscript{Log})} combinations
& \textbackslash{} \emph{F}\textsubscript{LOGIC} → 0.00 routes
8\%-base-rate failures to dispute; current ranks \#161/245 & −4.43pp \\
3 & 1 & Threshold optimization & Grid search over 62 (upper, lower)
pairs & Optimal band 0.45/0.40 is 6× tighter than current 0.60/0.30 &
−3.97pp \\
4 & 1 & Weight sensitivity & ±5/10/15/20\% perturbation on each weight &
\textbackslash{} \emph{w\textsubscript{f}} most sensitive (+4.56pp at
+20\%); stable within ±10\% & --- \\
5 & 1 & Cross-task LOO-CV & Leave-one-type-out across 5 task types &
37.65\% ± 0.68\% --- generalizes across domains & --- \\
6 & 2 & Piecewise + interaction grid & Staged grid over
(\emph{w},\textbackslash{} \emph{b}\textsubscript{crit},\textbackslash{}
\emph{d}\textsubscript{crit}, penalties,\textbackslash{}
\emph{w}\textsubscript{int}) & Optimal:\textbackslash{}
\emph{b}\textsubscript{crit}=0.10,\textbackslash{}
\emph{d}\textsubscript{crit}=0.05,\textbackslash{}
\emph{w}\textsubscript{int}=0.25 & −0.64pp \\
7 & 2 & Ablation: cliff vs interaction & Compare cliff-only,
interaction-only, combined & Cliffs contribute \textbf{−0.01pp} (inert);
interaction contributes \textbf{−0.63pp} & --- \\
8 & 2 & Eq2 sensitivity & ±20\% perturbation on 8 Eq2 parameters & All 4
piecewise parameters have \textbf{zero sensitivity} --- structurally
inert & --- \\
9 & 3 & 5-variable weight optimization & Grid search
with\textbackslash{} \emph{w\textsubscript{c}},\textbackslash{}
\emph{w\textsubscript{s}} added & Optimizer assigns minimum (0.05) to
both new variables; at Eq1 thresholds the 5-var weight change alone
produces 34.98\% (+1.17pp vs Eq1-opt). Improvement arrives only when
re-tuned thresholds (Exp 10) are applied. & +1.17pp at Eq1 thresholds \\
10 & 3 & 5-variable threshold optimization & Threshold grid for Eq3
scores & Optimal: upper=0.50, lower=0.45 --- tighter than Eq1's
0.45/0.40 & −1.00pp \\
11 & 3 & Variable ablation & Solo complexity, solo skill, both & Solo:
−0.86pp (complexity), −0.87pp (skill); combined: −0.55pp
(\textbf{subadditive} --- linear sum cannot capture the multiplicative
\emph{C}·\emph{S} ground-truth interaction) & --- \\
12 & 3 & 5-var cross-task LOO-CV & Leave-one-type-out on Eq3 & 32.37\% ±
0.36\% --- generalizes; confirms\textbackslash{} \textasciitilde33\%
ceiling is structural, not data-specific & --- \\
13 & 4 & Bayes-optimal baseline & Route using ground-truth \emph{p}
directly & Binary floor 22.52\%; three-tier floor 22.53\% --- any
formula ≤25\% is near-optimal & --- \\
14 & 4 & Multiplicative grid search & 2,646 Phase-A configs (9×7×7
exponent triples × 6 coarse threshold pairs) + 53 Phase-B threshold
refinements at best exponents & Optimal: (0.80, 0.35, 0.15), thresholds
0.40/0.35 & −10.35pp vs Eq1 \\
15 & 4 & Hybrid α-sweep & \emph{r} = α·Eq1 + (1−α)·Eq4 at 11 α values &
Monotonic: α=0.0 best (23.46\%), α=1.0 worst (35.07\%) --- every
increment of linear component strictly degrades routing & --- \\
16 & 4 & Multiplicative cross-task LOO-CV & Leave-one-type-out on Eq4 &
23.39\% ± 0.36\% --- best generalization of any run & --- \\
\end{longtable}
}

\emph{Ground-truth model validation.} Before any optimization, the
synthetic ground truth was validated against published literature
{[}1{]}{[}2{]}{[}3{]}. The model is calibrated --- not invented --- from
prior empirical studies:

{\def\LTcaptype{none} % do not increment counter
\begin{longtable}[]{@{}
  >{\raggedright\arraybackslash}p{(\linewidth - 6\tabcolsep) * \real{0.2500}}
  >{\raggedright\arraybackslash}p{(\linewidth - 6\tabcolsep) * \real{0.2500}}
  >{\raggedright\arraybackslash}p{(\linewidth - 6\tabcolsep) * \real{0.2500}}
  >{\raggedright\arraybackslash}p{(\linewidth - 6\tabcolsep) * \real{0.2500}}@{}}
\toprule\noalign{}
\begin{minipage}[b]{\linewidth}\raggedright
Check
\end{minipage} & \begin{minipage}[b]{\linewidth}\raggedright
Expected (from literature)
\end{minipage} & \begin{minipage}[b]{\linewidth}\raggedright
Observed (simulation)
\end{minipage} & \begin{minipage}[b]{\linewidth}\raggedright
Deviation
\end{minipage} \\
\midrule\noalign{}
\endhead
\bottomrule\noalign{}
\endlastfoot
LIVENESS class frequency & \textasciitilde45\% & 44.96\% & 0.04pp \\
RESOURCE class frequency & \textbackslash{} \textasciitilde35\% &
34.87\% & 0.13pp \\
LOGIC class frequency & \textasciitilde20\% & 20.17\% & 0.17pp \\
LIVENESS recovery at high resources & \textbackslash{}
\textasciitilde92\% & 71.4\% & Base rate × complexity/skill factors \\
RESOURCE recovery at high resources & \textasciitilde48\% & 37.9\% &
Same modulation \\
LOGIC recovery at high resources & \textbackslash{} \textasciitilde8\% &
5.4\% & Same modulation \\
Overall recovery rate & \textasciitilde50\% {[}3{]} & 36.8\% &
Complexity/skill factors pull rate below base \\
\end{longtable}
}

Class frequencies match literature to within 0.17pp. Recovery rates are
proportionally scaled down by the complexity and skill factors, which
model real-world subtask length and fallback competence variance ---
factors absent from the base-rate literature but necessary for realistic
routing simulation.

\emph{Formula-level comparison (Eq1 optimized vs Eq4 multiplicative).}
The confusion matrices reveal where the multiplicative formula wins:

{\def\LTcaptype{none} % do not increment counter
\begin{longtable}[]{@{}
  >{\raggedright\arraybackslash}p{(\linewidth - 8\tabcolsep) * \real{0.2000}}
  >{\raggedright\arraybackslash}p{(\linewidth - 8\tabcolsep) * \real{0.2000}}
  >{\raggedright\arraybackslash}p{(\linewidth - 8\tabcolsep) * \real{0.2000}}
  >{\raggedright\arraybackslash}p{(\linewidth - 8\tabcolsep) * \real{0.2000}}
  >{\raggedright\arraybackslash}p{(\linewidth - 8\tabcolsep) * \real{0.2000}}@{}}
\toprule\noalign{}
\begin{minipage}[b]{\linewidth}\raggedright
Routing Cell
\end{minipage} & \begin{minipage}[b]{\linewidth}\raggedright
Eq1 Linear
\end{minipage} & \begin{minipage}[b]{\linewidth}\raggedright
Eq4 Multiplicative
\end{minipage} & \begin{minipage}[b]{\linewidth}\raggedright
Bayes Optimal
\end{minipage} & \begin{minipage}[b]{\linewidth}\raggedright
Meaning
\end{minipage} \\
\midrule\noalign{}
\endhead
\bottomrule\noalign{}
\endlastfoot
FULL + Succeeded & 26.1\% & 21.8\% & 23.0\% & Correctly routed to full
recovery \\
\textbf{FULL + Failed} & \textbf{22.3\%} & \textbf{7.9\%} &
\textbf{9.2\%} & \textbf{False positive --- wasted recovery} \\
REDUCED + Succeeded & 3.3\% & 3.2\% & 0.4\% & Correctly routed to
reduced recovery \\
REDUCED + Failed & 3.8\% & 5.4\% & 0.1\% & False positive in reduced
tier \\
\textbf{DISPUTED + Succeeded} & \textbf{7.7\%} & \textbf{12.2\%} &
\textbf{13.0\%} & \textbf{False negative --- missed recovery} \\
DISPUTED + Failed & 36.7\% & 51.5\% & 53.3\% & Correctly disputed \\
\end{longtable}
}

The headline result: \textbf{FULL-tier false positives drop from 22.3\%
to 7.9\% --- a 65\% reduction}. The multiplicative formula is far more
selective about which tasks receive full recovery resources. The
trade-off is a rise in disputed-but-recoverable cases (7.7\% → 12.2\%),
which is the correct direction: a failed recovery wastes the fallback's
budget and time, while a disputed-recoverable task merely delays
resolution with the arbiter fee as overhead (see Section 6.6 for
economic cost). The Eq4 matrix is strikingly close to the Bayes-optimal
matrix, confirming the formula captures nearly all information
extractable from the three on-chain inputs.

\emph{Cross-task-type generalization (LOO-CV).} The per-task-type
leave-one-out results confirm the formula is not over-fit to any single
domain:

{\def\LTcaptype{none} % do not increment counter
\begin{longtable}[]{@{}
  >{\raggedright\arraybackslash}p{(\linewidth - 6\tabcolsep) * \real{0.2500}}
  >{\raggedright\arraybackslash}p{(\linewidth - 6\tabcolsep) * \real{0.2500}}
  >{\raggedright\arraybackslash}p{(\linewidth - 6\tabcolsep) * \real{0.2500}}
  >{\raggedright\arraybackslash}p{(\linewidth - 6\tabcolsep) * \real{0.2500}}@{}}
\toprule\noalign{}
\begin{minipage}[b]{\linewidth}\raggedright
Held-out Task Type
\end{minipage} & \begin{minipage}[b]{\linewidth}\raggedright
Run 1 (Eq1)
\end{minipage} & \begin{minipage}[b]{\linewidth}\raggedright
Run 4 (Eq4)
\end{minipage} & \begin{minipage}[b]{\linewidth}\raggedright
Eq4 Improvement
\end{minipage} \\
\midrule\noalign{}
\endhead
\bottomrule\noalign{}
\endlastfoot
\texttt{defi.price\_fetch} & 37.90\% & 23.96\% & −13.94pp \\
\texttt{defi.trade\_execute} & 38.23\% & 23.32\% & −14.91pp \\
\texttt{data.report\_generate} & 36.63\% & 22.83\% & −13.80pp \\
\texttt{governance.vote\_delegate} & 38.39\% & 23.39\% & −15.00pp \\
\texttt{compute.model\_inference} & 37.10\% & 23.44\% & −13.66pp \\
\textbf{Mean ± std} & \textbf{37.65\% ± 0.68\%} & \textbf{23.39\% ±
0.36\%} & \textbf{−14.26pp ± 0.54pp} \\
\end{longtable}
}

Both formulas generalize (std well below the 3pp threshold), but Eq4
generalizes twice as tightly (0.36\% vs 0.68\%).
\texttt{compute.model\_inference} is the worst case for Eq4 (23.44\%)
--- still 9pp below the Eq1 best case. The reported standard deviations
use the population formula (divisor \emph{N}, consistent with NumPy's
default in \texttt{results\_eq4.json}); the sample-std form (divisor
\emph{N}−1) yields 0.40\% for Eq4 and 0.76\% for Eq1. The
\emph{N}-divisor convention is retained here to match the raw simulation
output.

\emph{Ground-truth model specification.} The simulation uses base
recovery rates from the MAST taxonomy {[}1{]} and agent reliability
literature {[}2{]}{[}3{]}: LIVENESS 92\%, RESOURCE 48\%, LOGIC 8\%. The
recovery probability for any single task-failure event is:

\begin{verbatim}
p = base × σ(15·(B − 0.15)) × σ(20·(D − 0.10)) × 1/(1 + 0.02·n_remaining) × (0.4 + 0.6·skill)
\end{verbatim}

where σ is the logistic function, \emph{B} and \emph{D} are budget and
deadline remaining, \emph{n\_remaining} is the count of remaining
subtasks, and \emph{skill} ∈ {[}0, 1{]} is the fallback agent's drawn
skill score. The steeper deadline sigmoid (slope 20 vs budget's 15)
reflects the sharper empirical cliff around deadline exhaustion. Skill
is gated by the fallback's ERC-8004 reputation: higher-reputation
fallbacks sample from a distribution biased toward 1.0. The source of
truth is \texttt{simulation/recovery.py}; the stochastic event generator
and all 16 experiment scripts live in \texttt{simulation/}, with results
in \texttt{RESULTS.md} (Run 1), \texttt{RESULTS\_EQ2.md} (Run 2),
\texttt{RESULTS\_EQ3.md} (Run 3), \texttt{RESULTS\_EQ4.md} (Run 4).

\textbf{Remaining calibration work.} The simulation uses synthetic
failure distributions. As CAIRN accumulates real execution records, the
exponents (\emph{a}, \emph{b}, \emph{c}), class weights, and thresholds
should be re-validated against observed recovery outcomes via Bayesian
posterior inference (MCMC) and adjusted via governance. The staged
calibration roadmap: \textbf{(Stage 1)} Monte Carlo with synthetic data
--- \textbf{complete} (this section); \textbf{(Stage 2)} Bayesian
posterior update once testnet produces ≥ \emph{k} = 30 records per class
per task type\textbackslash{} (\textasciitilde{} 60 days post-launch at
10 tasks/day, per Section 5.3) --- \textbf{planned}; \textbf{(Stage 3)}
full MCMC re-calibration at mainnet with per-type posteriors ---
\textbf{planned}. At each stage the exponents are updated via governance
parameter change, not a code redeployment.

\textbf{Multi-agent recovery chains.} The current protocol supports one
fallback. Future versions could support multiple sequential fallbacks,
with each contributing checkpoints and earning proportional payment. The
mechanism design for chains longer than two agents requires additional
analysis.

\subsection{10.2 Protocol Extensions}\label{protocol-extensions}

\textbf{ERC standardization.} CAIRN is designed to become an Ethereum
standard. The specification in \href{./ERC-CAIRN.md}{ERC-CAIRN.md}
follows the EIP-1 format. Working title:
\texttt{ERC-CAIRN:\ Agent\ Failure\ and\ Recovery\ Standard}.

\textbf{CAIRN MCP Server.} Exposing checkpoint, recovery, and
intelligence query as MCP tools would enable any MCP-connected agent to
participate in CAIRN without framework-specific integration.

\textbf{Cross-chain support.} CAIRN is currently deployed on Base.
Cross-chain fallback (e.g., a Base task recovered by an Olas agent on
Gnosis) requires cross-chain messaging and is planned for a future
protocol version.

\textbf{Privacy-preserving intelligence.} Currently, all failure records
are public. Future versions may use zero-knowledge proofs to enable
agents to query failure patterns without revealing their specific
execution data.

\subsection{10.3 Limitations}\label{limitations}

CAIRN's current design makes explicit trade-offs. We state them here to
bound the claims made elsewhere in this paper.

\textbf{Checkpoint portability boundary.} CAIRN's full checkpoint
portability covers structured pipeline tasks (approximately 90\% of
current on-chain agent workloads). Reasoning-heavy tasks
(chain-of-thought, planning with backtracking) operate in degraded mode
where only output-level checkpoints are portable. See Section 4.1.1.

\textbf{Recovery score accuracy.} The multiplicative formula achieves
23.46\% misrouting against the synthetic ground truth --- 0.93pp from
the Bayes-optimal floor of 22.53\% (96\% of achievable improvement
captured). The ground truth is calibrated to published class frequencies
to within 0.17pp (LIVENESS/RESOURCE/LOGIC at 44.96\% / 34.87\% / 20.17\%
vs.~literature targets 45\% / 35\% / 20\%), but the recovery
\emph{rates} within each class have not been validated against empirical
recovery outcomes because no such dataset yet exists. The staged
calibration roadmap (Section 10.1) replaces the synthetic ground truth
with observed outcomes as testnet and mainnet data accumulate; parameter
updates happen via governance, not redeployment.

\textbf{Single-fallback architecture.} CAIRN supports one fallback
attempt per failure. If the fallback also fails, the task goes to
dispute. Multi-fallback chains are deferred to a future version.

\textbf{On-chain classification limits.} Failure classification is
high-confidence for LIVENESS (heartbeat miss) and RESOURCE
(budget/deadline exceeded) triggers, but LOGIC failures require agent
self-reporting or external verification --- a weaker signal. See the
LOGIC class definition in Section 3.1 and the \texttt{detectFailure}
permissionless-enforcement discussion in Section 3.3.

\textbf{L2 dependency.} CAIRN is economically viable on Base L2 but not
on Ethereum mainnet (gas costs would be 100-1000× higher). This creates
a dependency on the L2's sequencer availability and ordering guarantees.

\textbf{Arbiter recursion.} Arbiters are themselves agents reasoning
about reasoning errors. Proposition 3 (Section 7.5) establishes that
honest ruling is rational when the detection probability
\emph{p}\textsubscript{d} for incorrect rulings is high --- which holds
for LIVENESS and RESOURCE disputes (on-chain evidence) but not
necessarily for LOGIC disputes (``did the agent hallucinate?''), where
detection is itself a reasoning inference. CAIRN inherits the structural
limitation common to all on-chain arbitration protocols {[}19{]}: a
sufficiently adversarial LOGIC dispute can have a genuinely contestable
answer, and the protocol cannot fully economically deter collusion in
that regime. Mitigations (reputation gates, commit-reveal, multi-dispute
compounding) bound the worst case; a Kleros-style multi-juror appeals
layer is reserved for protocol v3.

\textbf{Arbiter depth.} Relatedly, the current arbiter mechanism is
single-tier with no appeals. High-value disputes may require additional
dispute resolution infrastructure in future versions.

\textbf{v1 / v2 specification gap.} The v1 testnet contract implements
the pre-calibration linear recovery formula (see Section 6.4 note), 15\%
arbiter stake, binary routing at threshold 0.30, a non-batched
\texttt{commitCheckpoint} signature in the MVP variant, and no on-chain
schema-hash enforcement. The v2 specification described throughout this
paper --- multiplicative formula, three-tier routing at 0.40/0.35, 20\%
arbiter stake, batched \texttt{commitCheckpointBatch}, PRBMath-based or
lookup-based fixed-point exponentiation, and per-checkpoint schema
validation --- is the target protocol that the simulation work
motivates. The migration path is governance-gated via the
\texttt{IRecoveryRouter} interface; it does not require a state-breaking
upgrade of deployed tasks. Gas figures (Section 6.5) are design-target
estimates for v2; the \texttt{forge\ test\ -\/-gas-report} against the
v2 reference implementation will be published alongside the v2 testnet
deployment. Readers evaluating this paper should treat
deployed-behaviour claims as pertaining to v2 unless explicitly
annotated ``v1.''

\begin{center}\rule{0.5\linewidth}{0.5pt}\end{center}

\section{11. References}\label{references}

{[}1{]} M. Cemri, M. Z. Pan, S. Yang, et al., ``Why Do Multi-Agent LLM
Systems Fail?'', \emph{NeurIPS 2025 Datasets and Benchmarks Track},
arXiv:2503.13657, 2025.

{[}2{]} S. Rabanser, S. Kapoor, P. Kirgis, K. Liu, S. Utpala, A.
Narayanan, ``Towards a Science of AI Agent Reliability'',
arXiv:2602.16666, February 2026.

{[}3{]} R. Lu, Y. Li, Y. Huo, ``Exploring Autonomous Agents: A Closer
Look at Why They Fail'', \emph{ASE 2025 NIER Track}, arXiv:2508.13143,
August 2025.

{[}4{]} ``Blockchain-Enhanced Incentive-Compatible Mechanisms for
Multi-Agent Reinforcement Learning Systems'', \emph{Nature Scientific
Reports}, November 2025. DOI: 10.1038/s41598-025-20247-8.

{[}5{]} Olas Network, ``Mech Marketplace'',
\url{https://olas.network/mech-marketplace}. Over 10 million
agent-to-agent transactions as of 2026; approximately 2,000 agents
deployed,\textbackslash{} \textasciitilde500 active daily.

{[}6{]} S. Alqithami, ``Autonomous Agents on Blockchains: Standards,
Execution Models, and Trust Boundaries'', arXiv:2601.04583, January
2026. Systematic survey of 317 publications identifying missing
interface layers and verifiable policy enforcement as key gaps.

{[}7{]} ISO/IEC TR 5469:2024, ``Artificial Intelligence --- Functional
Safety and AI Systems'', International Organization for Standardization,
2024.

{[}8{]} K. M. Chandy and L. Lamport, ``Distributed Snapshots:
Determining Global States of Distributed Systems'', \emph{ACM
Transactions on Computer Systems}, 3(1):63-75, 1985.

{[}9{]} V. Buterin, D. Ryan, et al., ``Ethereum Proof-of-Stake Consensus
Specifications'', Ethereum Foundation, 2020-2026.
\url{https://github.com/ethereum/consensus-specs}. Staking/slashing
mechanism securing over \$100 billion in staked value across 1M+
validators as of early 2026.

{[}10{]} ``AI Agents Meet Blockchain: A Survey'', \emph{MDPI Future
Internet}, 17(2):57, February 2025. Introduces the Proof-of-Thought
concept.

{[}11{]} IETF Draft, ``Task-Oriented Multi-Agent Recovery Framework for
High-Reliability in Converged Mobile Networks'',
\texttt{draft-yue-anima-agent-recovery-networks-00}, 2026.

{[}12{]} Google, ``Agent2Agent Protocol (A2A)'', donated to Linux
Foundation June 2025; v0.3 released July 2025.
\url{https://github.com/a2aproject/A2A}

{[}13{]} Anthropic, ``Model Context Protocol (MCP)'', 2024-2026; donated
to Linux Foundation Agentic AI Foundation December 2025.
\url{https://modelcontextprotocol.io}

{[}14{]} M. De Rossi, D. Crapis, J. Ellis, E. Reppel, ``ERC-8004:
Trustless Agents Standard.'' EIP in Draft status; mainnet-deployed since
January 29, 2026. \textasciitilde49,000 registered agents across 30+ EVM
chains as of February 2026. EIP:
\url{https://eips.ethereum.org/EIPS/eip-8004}

{[}15{]} D. Crapis, B. Lim, T. Weixiong, C. Zuhwa, ``ERC-8183: Agentic
Commerce Standard.'' EIP in Draft status, created February 25, 2026.
Agent Commerce Protocol (ACP) deployed on Arbitrum via Virtuals
Protocol. EIP: \url{https://eips.ethereum.org/EIPS/eip-8183}

{[}16{]} Biconomy, ``ERC-8211: Smart Batching --- Runtime-Resolved
Parameters and Predicate-Gated Execution for Smart Accounts.'' Proposed
April 2026; discussion draft on Ethereum Magicians (no canonical EIP
page at time of writing). Discussion:
\url{https://ethereum-magicians.org/t/erc-8211-smart-batching/28135}

{[}17{]} R. McPeck, D. Finlay, R. Dawson, D. Chiang, ``ERC-7710: Smart
Contract Delegation.'' EIP in Draft status, created May 20, 2024. Used
by the MetaMask Delegation Toolkit. EIP:
\url{https://eips.ethereum.org/EIPS/eip-7710}

{[}18{]} CAIRN Recovery Score Calibration Simulation, April 2026. Monte
Carlo validation across 100,000 synthetic task-failure events per run, 4
formula structures (linear, piecewise + interaction, 5-variable linear,
multiplicative), 16 experiments (Exp 1-5
weight/class/threshold/sensitivity/LOO-CV for Eq1; Exp 6-8 for Eq2; Exp
9-12 for Eq3; Exp 13-16 Bayes-optimal baseline, multiplicative grid,
hybrid α-sweep, and cross-task LOO-CV for Eq4). Run 1: 362 grid points
(55 weight + 245 class-weight + 62 threshold), runtime\textbackslash{}
\textasciitilde3 seconds. Run 2: staged grid over 8 Eq2 parameters,
runtime \textasciitilde31 seconds. Run 3: staged grid over 5 linear
weights + thresholds. Run 4: 2,646 multiplicative-formula grid points +
hybrid α-sweep, runtime\textbackslash{} \textasciitilde14 seconds.
Reproducible: \texttt{python3\ -m\ simulation.run} (Run 1),
\texttt{run\_eq2} (Run 2), \texttt{run\_eq3} (Run 3), \texttt{run\_eq4}
(Run 4); seed=42, deterministic on any NumPy ≥1.20 installation. Source:
\texttt{simulation/} in the CAIRN repository. Results:
\texttt{simulation/RESULTS.md}, \texttt{RESULTS\_EQ2.md},
\texttt{RESULTS\_EQ3.md}, \texttt{RESULTS\_EQ4.md}. Figures:
\texttt{simulation/figures/fig1} through \texttt{fig16}.

{[}19{]} C. Lesaege, F. Ast, W. George, ``Kleros: A Decentralized Court
System for the Internet'', Kleros Yellowpaper, v1.0.7, 2019.
\url{https://kleros.io/yellowpaper.pdf}. Pioneers the
bonded-juror-with-Schelling-point design for on-chain dispute
resolution, including multi-round appeals with quadratically increasing
juror pools. Cited in this paper as the canonical prior work on the
arbiter-recursion problem (Section 4.5, Section 7.5, Section 10.3).

\begin{center}\rule{0.5\linewidth}{0.5pt}\end{center}

\emph{CAIRN --- Agent Failure and Recovery Protocol} \emph{Whitepaper
v2.0 --- April 2026} \emph{Agents learn together.}

\end{document}
